\documentclass[11pt]{article}
\usepackage[margin=1.1in]{geometry}
\usepackage{amsmath,amssymb,amsthm,mathtools}
\usepackage{enumitem}
\usepackage{hyperref}
\usepackage{microtype}
\usepackage{booktabs}
\usepackage{mathrsfs}
\hypersetup{colorlinks=true,linkcolor=blue,citecolor=blue,urlcolor=blue,hypertexnames=false,pdftitle={Endpoint Short Cycles},pdfauthor={Christopher D. Long}}
\newtheorem{theorem}{Theorem}[section]
\newtheorem{proposition}[theorem]{Proposition}
\newtheorem{lemma}[theorem]{Lemma}
\newtheorem{corollary}[theorem]{Corollary}
\newtheorem{conjecture}[theorem]{Conjecture}
\theoremstyle{remark}
\newtheorem{remark}[theorem]{Remark}
\newtheorem{warning}[theorem]{Warning}
\newtheorem{definition}[theorem]{Definition}
\newcommand{\Pp}{\mathbb P}
\newcommand{\Ee}{\mathbb E}
\newcommand{\Rr}{\mathbb R}
\newcommand{\Nn}{\mathbb N}
\newcommand{\Zz}{\mathbb Z}
\newcommand{\one}{\mathbf 1}
\newcommand{\TV}{\mathrm{TV}}
\newcommand{\Pois}{\operatorname{Pois}}
\newcommand{\Exp}{\operatorname{Exp}}
\newcommand{\PPP}{\operatorname{PPP}}
\newcommand{\Var}{\operatorname{Var}}
\newcommand{\Cov}{\operatorname{Cov}}
\newcommand{\diag}{\operatorname{diag}}
\newcommand{\dd}{\,d}
\newcommand{\eps}{\varepsilon}
\newcommand{\calD}{\mathcal D}
\newcommand{\calC}{\mathcal C}
\newcommand{\calE}{\mathcal E}
\newcommand{\calG}{\mathcal G}
\title{Endpoint Short Cycles}
\author{Christopher D. Long}
\date{Mechanical extraction --- proof audit pending}
\begin{document}\maketitle
\begin{abstract}
For uniformly endpoint-regularly varying Luce arrays, this paper identifies the logarithmically diverging mean of every fixed cycle length and proves joint fixed-order factorial-moment asymptotics. The coefficients are diagonal convolution traces of the multiplicative endpoint kernel. The proof combines source-to-target endpoint saddle estimates, a global scheduled-deletion cycle envelope, continuum and discrete logarithmic trace limits, and factorial-moment summation.
\end{abstract}
\tableofcontents
\section{Status and provenance}
This source is a mechanical extraction from the frozen 139-page Luce manuscript at repository commit \texttt{6178a99a8f2ee7b6416edd026e7780b9b8fa9fa0}. The mathematical blocks between the archival sentinels below are copied byte-for-byte from the frozen source. Only the standalone preamble, abstract, section wrappers, provenance notice, and cited bibliography entries have been added.

The extraction preserves every theorem hypothesis, identity, limiting statement, topology, quantifier, and proof step in the selected theorem package. It does not audit, repair, strengthen, weaken, or modernize the mathematics. Its status remains a working mechanical extraction until a separate adversarial audit is completed.

\section{Exponential and Gumbel input}

% BEGIN ARCHIVED BLOCK: exponential-race-input
\subsection{The exponential race}
Let \(\theta_1,\ldots,\theta_n>0\), and let \(X_i\sim\Exp(\theta_i)\) be independent.  For distinct labels \(i_1,\ldots,i_n\),
\[
\Pp(X_{i_1}<\cdots<X_{i_n})
= \prod_{r=1}^n
 \frac{\theta_{i_r}}{\theta_{i_r}+\theta_{i_{r+1}}+\cdots+\theta_{i_n}}.
\]
This is the standard exponential race identity and follows from the memoryless property.


% END ARCHIVED BLOCK: exponential-race-input
% BEGIN ARCHIVED BLOCK: gumbel-threshold-input
\subsection{Gumbel priorities}\label{subsec:gumbel-priorities}
A standard Gumbel random variable has distribution
\[
        \Pp(G\le y)=\exp(-e^{-y}),
        \qquad y\in\Rr .
\]
If \(E_i\sim\Exp(1)\) are independent, then \(-\log E_i\) are independent standard Gumbels.  Since
\[
        X_i\stackrel{d}=E_i/\theta_i,
\]
the exponential race may equivalently be generated by the shifted priorities
\[
        Y_i=-\log X_i=\log\theta_i+G_i,
        \qquad 1\le i\le n,
\]
ordered in decreasing order.  Thus
\[
        \sigma(i)=1+\#\{j:Y_j>Y_i\}.
\]
For later reference we record the exact threshold formulas.  The density and upper-tail function of \(Y_i\) are
\[
        f_i(y)=\theta_i e^{-y}\exp(-\theta_i e^{-y}),
        \qquad
        \Pp(Y_i>y)=1-\exp(-\theta_i e^{-y}).
\]

\begin{lemma}[Gumbel threshold identities]\label{lem:gumbel-threshold-identities}
For a single specified label \(p\),
\[
        \Pp(\sigma(p)=q)
        =\int_{-\infty}^{\infty}
        f_p(y)\,
        \Pp\left(\sum_{j\ne p}\one_{\{Y_j>y\}}=q-1\right)\dd y .
\]
More generally, let \(p_1,\ldots,p_\ell\) be distinct labels and let
\(q_1<\cdots<q_\ell\) be distinct target positions.  Put
\(P=\{p_1,\ldots,p_\ell\}\) and
\[
        S_a(y)=\sum_{j\notin P}\one_{\{Y_j>y_a\}},
        \qquad 1\le a\le\ell .
\]
Then
\[
\begin{aligned}
&\Pp(\sigma(p_a)=q_a,\ 1\le a\le\ell) \\
&\quad=\int_{y_1>\cdots>y_\ell}
        \prod_{a=1}^\ell f_{p_a}(y_a)\,
        \Pp(S_a(y)=q_a-a,\ 1\le a\le\ell)
        \dd y_1\cdots\dd y_\ell .
\end{aligned}
\]
\end{lemma}

\begin{proof}
Condition on the specified priorities.  For one label, the event
\(\sigma(p)=q\) is exactly the event that precisely \(q-1\) nonspecified
priorities exceed the observed value of \(Y_p\).  This gives the first
identity.

For the multi-point identity, the ordering of the target positions forces the
specified priorities to satisfy \(Y_{p_1}>\cdots>Y_{p_\ell}\).  Conditional on
\(Y_{p_a}=y_a\) in this chamber, exactly \(a-1\) specified priorities are above
\(y_a\).  Hence \(\sigma(p_a)=q_a\) is equivalent to
\(S_a(y)=q_a-a\), simultaneously for \(1\le a\le\ell\).  Integrating against the
independent specified-priority densities gives the formula.
\end{proof}

% END ARCHIVED BLOCK: gumbel-threshold-input
% BEGIN ARCHIVED BLOCK: analytic-input
\subsection{Analytic input}
We use standard estimates for sums of independent Bernoulli variables.  They are stated explicitly to make clear what is invoked.

\begin{lemma}[Point probability and Bernstein bounds]\label{lem:bernstein-point}
Let \(Y=\sum_{j=1}^N \xi_j\), where the \(\xi_j\) are independent, not necessarily identically distributed, Bernoulli variables, and set \(\mu=\Ee Y\), \(v=\Var Y\).  There is an absolute constant \(C\) such that
\[
        \sup_m \Pp(Y=m)\le \frac{C}{\sqrt{1+v}}.
\]
Moreover, there is an absolute constant \(c>0\) such that, for every \(u\ge0\),
\[
        \Pp(|Y-\mu|\ge u)\le
        2\exp\left[-c\min\left(\frac{u^2}{v+1},u\right)\right].
\]
\end{lemma}

\begin{proof}
The point bound is the classical concentration-function estimate for sums of Bernoulli variables.  It follows from Fourier inversion and the standard inequality
\[
        |1-p+pe^{it}|
        \le \exp\left[-\frac{2}{\pi^2}p(1-p)\operatorname{dist}(t,2\pi\mathbb Z)^2\right],
\]
where \(\operatorname{dist}(t,2\pi\mathbb Z)\) denotes distance from \(t\) to the nearest multiple of \(2\pi\).  Bernstein's inequality gives the second bound for bounded independent summands.  See Petrov \cite[Chapter VII]{Petrov1975} for both estimates.
\end{proof}

\begin{lemma}[Uniform lattice local CLT with a tilted local bound]
\label{lem:llt}
Fix a dimension $d$ and a support bound $B<\infty$.  Let
\[
        Y_{n,m}=\sum_{j=1}^{N_n}\zeta_{n,m,j},
        \qquad m\in\mathcal M_n,
\]
be sums of independent $\mathbb Z^d$-valued random vectors satisfying
$\|\zeta_{n,m,j}\|\le B$.  Write
\[
        \mu_{n,m}=\Ee Y_{n,m},
        \qquad \Sigma_{n,m}=\Var(Y_{n,m}),
\]
and suppose that, uniformly in $m$,
\[
        c s_n I_d\le\Sigma_{n,m}\le C s_n I_d,
        \qquad s_n\longrightarrow\infty.
\]
For $\lambda\in\mathbb R^d$, let $\Pp_\lambda$ denote the product exponential
tilt
\[
 \Pp_\lambda(\zeta_{n,m,j}=x)
 =\frac{e^{\lambda\cdot x}\Pp(\zeta_{n,m,j}=x)}
 {\Ee e^{\lambda\cdot\zeta_{n,m,j}}}.
\]
Assume that there is $\lambda_0>0$ for which the following two properties hold
uniformly for $\|\lambda\|\le\lambda_0$:
\begin{enumerate}[label=(\roman*)]
\item the tilted covariance matrices satisfy
\[
        c_0s_nI_d\le\Var_\lambda(Y_{n,m})\le C_0s_nI_d;
\]
\item for every $\varepsilon>0$ there are $\kappa_\varepsilon>0$ and
$n_\varepsilon$ such that
\[
 \left|\prod_{j=1}^{N_n}
 \Ee_\lambda e^{iu\cdot\zeta_{n,m,j}}\right|
 \le e^{-\kappa_\varepsilon s_n}
\]
whenever $n\ge n_\varepsilon$ and
$u\in[-\pi,\pi]^d$ has distance at least $\varepsilon$ from
$2\pi\mathbb Z^d$.
\end{enumerate}
Then, uniformly in $m$ and $z\in\mathbb Z^d$,
\begin{align*}
 \Pp(Y_{n,m}=z)
 &=(2\pi)^{-d/2}(\det\Sigma_{n,m})^{-1/2}\\
 &\quad\times
 \exp\left[-\frac12(z-\mu_{n,m})^T
 \Sigma_{n,m}^{-1}(z-\mu_{n,m})\right]
 +o(s_n^{-d/2}).
\end{align*}
Moreover,
\[
        \sup_z\Pp(Y_{n,m}=z)\le C_1s_n^{-d/2}.
\]
For every fixed $K<\infty$ and $\gamma<1/2$, there are $c_1,C_2>0$ such
that, uniformly when
\[
        \|z-\mu_{n,m}\|\le K s_n^{1/2+\gamma},
\]
one has
\[
 \Pp(Y_{n,m}=z)
 \le C_2s_n^{-d/2}
 \exp\left[-c_1\frac{\|z-\mu_{n,m}\|^2}{s_n}\right].
\]
All constants depend only on the fixed dimension, the support bound, and the
uniform constants in the hypotheses.
\end{lemma}

\begin{proof}
We suppress $n,m$ and write $Y,\mu,\Sigma,s$.  Let
$X_j=\zeta_j-\Ee\zeta_j$ and
\[
        \Psi(u)=\Ee e^{iu\cdot(Y-\mu)}.
\]
Since $\|X_j\|\le2B$ and
$\sum_j\Ee\|X_j\|^2=\operatorname{tr}\Sigma=O(s)$,
\[
        \sum_j\Ee\|X_j\|^3=O(s).
\]
Consequently, uniformly for $\|u\|$ sufficiently small,
\[
        \log\Psi(u)=-\frac12u^T\Sigma u+O(s\|u\|^3).
        \tag{2.1}
\]
A symmetrization argument also gives, on a fixed neighborhood of the origin,
\[
        |\Psi(u)|\le e^{-cs\|u\|^2}.
        \tag{2.2}
\]
Indeed, if $X_j'$ is an independent copy of $X_j$, then
$1-|\Ee e^{iu\cdot X_j}|^2=
\Ee[1-\cos(u\cdot(X_j-X_j'))]$; for small $u$ the right-hand side is
bounded below by a constant multiple of $\Var(u\cdot X_j)$, and multiplication
over $j$ yields (2.2).

Fourier inversion on $[-\pi,\pi]^d$ now proves the local expansion.  Choose
$L_s\to\infty$ with $L_s^3/\sqrt{s}\to0$ and split the integral into
\[
 \|u\|\le L_s/\sqrt{s},\qquad
 L_s/\sqrt{s}<\|u\|\le\varepsilon,
 \qquad \|u\|>\varepsilon.
\]
On the first region, (2.1), followed by the change of variables
$v=\Sigma^{1/2}u$, gives the Gaussian Fourier integral with an
$o(s^{-d/2})$ error.  The estimate (2.2) makes the second region
$o(s^{-d/2})$, and tilted aperiodicity with $\lambda=0$ makes the third
region exponentially small.  The oscillatory factor
$e^{-iu\cdot(z-\mu)}$ has modulus one, so every error is uniform in $z$.
Taking absolute values in the same inversion formula proves the concentration
bound.  This is the usual fixed-dimensional Fourier proof of the lattice local
central limit theorem; compare Petrov \cite[Chapter~VII]{Petrov1975}.

It remains to prove the Gaussian local upper bound.  Let
$h=z-\mu$ and define the centered cumulant generating function
\[
        K(\lambda)=\log\Ee e^{\lambda\cdot(Y-\mu)}.
\]
For $\|\lambda\|\le\lambda_0$, its Hessian is
$\nabla^2K(\lambda)=\Var_\lambda(Y)$ and is therefore between
$c_0sI_d$ and $C_0sI_d$.  If $\|h\|=o(s)$, strong convexity shows that the
minimum of $K(\lambda)-h\cdot\lambda$ over the ball
$\|\lambda\|\le\lambda_0/2$ is attained in its interior for all large $n$.
Thus there is a unique $\lambda_h$ in that ball satisfying
\[
        \nabla K(\lambda_h)=h,
        \qquad \|\lambda_h\|\le C\|h\|/s.
\]
The associated Legendre transform
$I(h)=\lambda_h\cdot h-K(\lambda_h)$ satisfies
\[
        I(h)\ge c\|h\|^2/s.
\]
For example, the last inequality follows by testing the variational formula
$I(h)=\sup_\lambda\{\lambda\cdot h-K(\lambda)\}$ at a sufficiently small
multiple of $h/s$ and using the tilted covariance upper bound.

Under $\Pp_{\lambda_h}$, the mean of $Y$ is $z$.  Repeating the absolute-value
Fourier estimate above, now using the two tilted hypotheses, gives
\[
        \sup_w\Pp_{\lambda_h}(Y=w)\le Cs^{-d/2}.
\]
The change-of-measure identity therefore yields
\[
 \Pp(Y=z)=e^{-I(h)}\Pp_{\lambda_h}(Y=z)
 \le Cs^{-d/2}e^{-c\|h\|^2/s}.
\]
The stated displacement range is $o(s)$ because $\gamma<1/2$, completing the
proof.
\end{proof}



% END ARCHIVED BLOCK: analytic-input

\section{Shared anisotropic local-limit toolkit}

% BEGIN ARCHIVED BLOCK: shared-anisotropic-toolkit
\begin{lemma}[Coordinatewise Fourier concentration]\label{lem:coord-fourier}
Let
\[
        W=\sum_{\nu\in\mathcal I}W_\nu
\]
be a sum of independent \(\mathbb Z^r\)-valued random vectors, each supported in a fixed cube \([-M_0,M_0]^r\).  Let
\(A_1,\ldots,A_r\ge1\).  Suppose that, for every coordinate \(j\), there is a set
\(\mathcal I_j\subset\mathcal I\) with
\(|\mathcal I_j|\ge c_0A_j\), and that the sets
\(\mathcal I_1,\ldots,\mathcal I_r\) are disjoint.  Suppose also that, for every
\(\nu\in\mathcal I_j\), the law of \(W_\nu\) gives mass at least \(c_0\) to two atoms whose difference is \(e_j\), the \(j\)-th coordinate vector.  Then there are constants
\(c,C>0\), depending only on \(r,M_0,c_0\), such that the characteristic function
\(\phi_W(u)=\Ee e^{iu\cdot W}\) satisfies
\[
        |\phi_W(u)|
        \le \exp\left[-c\sum_{j=1}^r
        A_j\operatorname{dist}(u_j,2\pi\mathbb Z)^2\right],
        \qquad u\in[-\pi,\pi]^r.
\]
Consequently,
\[
        \sup_z\Pp(W=z)\le C\prod_{j=1}^r A_j^{-1/2}.
\]
The same conclusions hold uniformly after any exponential tilt
\[
        \Pp_\lambda(W_\nu=x)
        =\frac{e^{\lambda\cdot x}\Pp(W_\nu=x)}
        {\Ee e^{\lambda\cdot W_\nu}},
        \qquad \|\lambda\|_\infty\le \Lambda,
\]
with constants depending additionally on \(\Lambda\).
\end{lemma}

\begin{proof}
For a single summand \(Y\), write \(p_x=\Pp(Y=x)\).  If two atoms
\(x\) and \(x+e_j\) both have mass at least \(c_0\), then
\[
\begin{aligned}
        1-|\Ee e^{iu\cdot Y}|^2
        &=\frac12\sum_{y,z}p_yp_z
        |e^{iu\cdot y}-e^{iu\cdot z}|^2  \\
        &\ge c_0^2 |1-e^{iu_j}|^2
        \ge c\,\operatorname{dist}(u_j,2\pi\mathbb Z)^2 .
\end{aligned}
\]
After decreasing \(c\), this gives
\[
        |\Ee e^{iu\cdot Y}|
        \le \exp\{-c\operatorname{dist}(u_j,2\pi\mathbb Z)^2\}.
\]
Multiplying this estimate over the disjoint blocks \(\mathcal I_j\) gives the displayed characteristic-function bound.  Fourier inversion then gives
\[
\begin{aligned}
        \Pp(W=z)
        &\le (2\pi)^{-r}\int_{[-\pi,\pi]^r}
        \exp\left[-c\sum_{j=1}^r
        A_j\operatorname{dist}(u_j,2\pi\mathbb Z)^2\right]\dd u \\
        &\le C\prod_{j=1}^r A_j^{-1/2}.
\end{aligned}
\]
Under a tilt with \(\|\lambda\|_\infty\le\Lambda\), all probabilities of atoms in the fixed support are multiplied by factors bounded between positive constants depending only on \(M_0,r,\Lambda\).  Hence the two distinguished atom probabilities remain bounded below, with a possibly smaller constant, and the same proof applies.
\end{proof}

\begin{lemma}[Weighted chain comparison]\label{lem:weighted-chain-comparison}
Fix $r\ge1$ and $\eta>0$.  Let
$0<A_1<\cdots<A_r$ satisfy $A_{j+1}/A_j\ge1+\eta$, and put
$u_{r+1}=0$.  Then
\[
        \sum_{j=1}^r A_j u_j^2
        \le C_{r,\eta}\sum_{j=1}^r A_j(u_j-u_{j+1})^2 .
\]
Equivalently, if $d_j=H_j-H_{j-1}$ with $H_0=0$, then
\[
        \sum_{j=1}^r \frac{d_j^2}{A_j}
        \asymp_{r,\eta}
        \sum_{j=1}^r \frac{H_j^2}{A_j} .
\]
\end{lemma}

\begin{proof}
Write $\Delta_j=u_j-u_{j+1}$.  Then $u_j=\sum_{k=j}^r\Delta_k$.  By Cauchy's inequality with the geometrically decaying weights $A_j/A_k\le(1+\eta)^{-(k-j)}$,
\[
        A_j u_j^2
        \le C_{r,\eta}\sum_{k=j}^r A_k\Delta_k^2 .
\]
Summing over $j$ gives the first display.  The second follows by applying the first inequality to the inverse triangular map and its reverse; since $r$ is fixed and the ratios are bounded below, both triangular maps have uniformly bounded operator norms in the weighted norms.
\end{proof}

\begin{lemma}[Truncated triangular comparison]\label{lem:truncated-triangular-comparison}
Fix $r\ge1$, $\eta>0$, and $M<\infty$.  Let
$0<A_1<\cdots<A_r$ satisfy $A_{j+1}/A_j\ge1+\eta$.  If
$d_j=H_j-H_{j-1}$, $H_0=0$, and $|d_j|\le M A_j$ for every $j$, then
\[
        \sum_{j=1}^r
        \min\left\{\frac{d_j^2}{A_j},A_j\right\}
        \ge c_{r,\eta,M}
        \sum_{j=1}^r
        \min\left\{\frac{H_j^2}{A_j},A_j\right\}.
\]
\end{lemma}

\begin{proof}
Write $x_j=d_j/A_j$.  Then $|x_j|\le M$ and
\[
        \frac{H_j}{A_j}=\sum_{k\le j}\frac{A_k}{A_j}x_k,
        \qquad
        \frac{A_k}{A_j}\le (1+\eta)^{-(j-k)} .
\]
Since $\min\{y^2,1\}\le y^2$, Cauchy's inequality and the geometric
bound on the ratios give, for each fixed $j$,
\[
\begin{aligned}
        A_j\min\left\{\frac{H_j^2}{A_j^2},1\right\}
        &\le A_j\left(\sum_{k\le j}\frac{A_k}{A_j}x_k\right)^2  \\
        &\le C_{r,\eta} A_j\sum_{k\le j}\frac{A_k}{A_j}x_k^2
         = C_{r,\eta}\sum_{k\le j} A_k x_k^2 .
\end{aligned}
\]
Because $|x_k|\le M$, we have
$x_k^2\le C_M\min\{x_k^2,1\}$.  Summing the preceding inequality over
$j$ and using that $r$ is fixed gives
\[
        \sum_{j=1}^r A_j
        \min\left\{\frac{H_j^2}{A_j^2},1\right\}
        \le C_{r,\eta,M}
        \sum_{k=1}^r A_k\min\{x_k^2,1\}.
\]
This is exactly the claimed inequality after rewriting
$A_j\min\{H_j^2/A_j^2,1\}=\min\{H_j^2/A_j,A_j\}$ and
$A_k\min\{x_k^2,1\}=\min\{d_k^2/A_k,A_k\}$.
\end{proof}

\begin{lemma}[Boundary-aware truncated chain comparison]
\label{lem:boundary-truncated-chain}
Fix \(r\ge1\) and \(M<\infty\).  Let
\(0<A_1<\cdots<A_r\), let \(B_-\) and \(B_+\) be disjoint subsets of
\(\{1,\ldots,r\}\), and put
\[
        G=\{1,\ldots,r\}\setminus(B_-\cup B_+).
\]
Let \(H_0=0\), \(d_j=H_j-H_{j-1}\), and
\[
        \phi_j(x)=\min\left\{\frac{x^2}{A_j},A_j\right\}.
\]
Assume
\[
        |d_j|\le M A_j\quad(j\in G),
        \qquad
        |H_j|\le M A_j\quad(j\in B_+).
\]
Then
\[
\begin{aligned}
&\sum_{j\in G}\phi_j(d_j)
 +\sum_{j\in B_-}|H_j|
 +\sum_{j\in B_+}A_j\\
&\qquad\ge c_{r,M}
 \sum_{j\in G}\phi_j(H_j).
\end{aligned}
\]
\end{lemma}

\begin{proof}
Decompose \(G\) into maximal consecutive runs.  Let
\(\{\ell,\ell+1,\ldots,u\}\) be one such run and put
\(b=\ell-1\).  For every \(j\) in the run,
\[
        H_j=H_b+\sum_{k=\ell}^j d_k,
\]
where \(H_b=H_0=0\) if \(\ell=1\).  For a fixed positive number \(A\) and
at most \(r+1\) real numbers \(x_i\),
\[
        \min\left\{\frac{(\sum_i x_i)^2}{A},A\right\}
        \le C_r\sum_i
        \min\left\{\frac{x_i^2}{A},A\right\}.
        \tag{7.1}
\]
Indeed, if one term on the right equals \(A\), the claim is immediate; otherwise
it follows from Cauchy's inequality.

If \(k\le j\) belongs to the run, then \(A_k\le A_j\) and
\(|d_k|\le M A_k\).  Hence
\[
        \phi_j(d_k)\le C_M\phi_k(d_k).
        \tag{7.2}
\]
To see this, the assertion is immediate when \(|d_k|\le A_k\); when
\(|d_k|>A_k\), one has
\(\phi_k(d_k)=A_k\) and
\(\phi_j(d_k)\le d_k^2/A_j\le M^2A_k\).

It remains to control the left boundary.  If \(b\in B_-\), then
\[
        \phi_j(H_b)\le |H_b|.
\]
If \(b\in B_+\), then \(|H_b|\le M A_b\) and \(A_b\le A_j\), so
\[
        \phi_j(H_b)\le M^2 A_b.
\]
Combining these estimates with (7.1) and (7.2), and then summing first over
\(j\) in the run and finally over the finitely many runs, proves the lemma.
\end{proof}

\begin{lemma}[Edge Fourier concentration]\label{lem:edge-fourier}
Let $e_1,\ldots,e_r$ be the coordinate vectors in $\mathbb Z^r$ and put $e_{r+1}=0$.  Let
\[
        W=\sum_{\nu\in\mathcal I}W_\nu
\]
be a sum of independent random vectors supported on
$\{e_1,\ldots,e_r,0\}$.  Let $A_1,\ldots,A_r\ge1$ satisfy
$A_{j+1}/A_j\ge1+\eta$ for $j<r$.  Suppose that, for every edge $j$, there is a set
$\mathcal I_j\subset\mathcal I$ with $|\mathcal I_j|\ge c_0A_j$, the sets $\mathcal I_j$ are disjoint, and for every $\nu\in\mathcal I_j$ the two atoms $e_j$ and $e_{j+1}$ each have probability at least $c_0$.  Then
\[
        |\Ee e^{iu\cdot W}|
        \le \exp\left[-c\sum_{j=1}^r
        A_j\operatorname{dist}(u_j-u_{j+1},2\pi\mathbb Z)^2\right],
        \qquad u_{r+1}=0,
\]
and consequently
\[
        \sup_z\Pp(W=z)\le C\prod_{j=1}^r A_j^{-1/2}.
\]
The same conclusions hold uniformly after any exponential tilt
$\|\lambda\|_\infty\le\Lambda$, with constants depending additionally on
$\Lambda$.
\end{lemma}

\begin{proof}
If a summand $Y$ gives mass at least $c_0$ to $e_j$ and to $e_{j+1}$, then
\[
        1-|\Ee e^{iu\cdot Y}|^2
        \ge c_0^2|e^{iu_j}-e^{iu_{j+1}}|^2
        \ge c\operatorname{dist}(u_j-u_{j+1},2\pi\mathbb Z)^2 .
\]
As in Lemma \ref{lem:coord-fourier}, this implies
\[
        |\Ee e^{iu\cdot Y}|
        \le \exp\{-c\operatorname{dist}(u_j-u_{j+1},2\pi\mathbb Z)^2\}.
\]
Multiplying over the disjoint edge blocks gives the displayed Fourier bound.
For the concentration estimate, Fourier inversion and Lemma
\ref{lem:weighted-chain-comparison} give
\[
\begin{aligned}
        \Pp(W=z)
        &\le (2\pi)^{-r}\int_{[-\pi,\pi]^r}
        \exp\left[-c\sum_{j=1}^r A_j
        \operatorname{dist}(u_j-u_{j+1},2\pi\mathbb Z)^2\right]\dd u \\
        &\le C\prod_{j=1}^r A_j^{-1/2}.
\end{aligned}
\]
Bounded tilts only multiply the finitely many atom probabilities by bounded factors, so the edge-block lower bounds persist after changing constants.
\end{proof}

\begin{lemma}[Forward-edge projected concentration]\label{lem:projected-forward-edge-concentration}
Let $G\ne\emptyset$ be a subset of $\{1,\ldots,r\}$ and let
$0<A_1<\cdots<A_r$ satisfy $A_{j+1}/A_j\ge1+\eta$.
Let $W$ be a sum of independent random vectors supported on
$\{e_j:j\in G\}\cup\{0\}$ after projection onto the coordinates in $G$.
Suppose that for each $j\in G$ there is an index
$\kappa(j)\in (G\cap\{j+1,\ldots,r\})\cup\{0\}$ and a set of at
least $c_0A_j$ summands, disjoint as $j$ varies, for which the two
projected atoms $e_j$ and $e_{\kappa(j)}$ have probabilities at least
$c_0$; here $e_0=0$.  Then, uniformly in $z_G$,
\[
        \Pp(W_G=z_G)\le C\prod_{j\in G} A_j^{-1/2}.
\]
The same estimate holds after any exponential tilt whose coordinates are
bounded in absolute value by a fixed constant.
\end{lemma}

\begin{proof}
For a summand in the block assigned to $j$, the same two-atom argument used
in Lemma \ref{lem:edge-fourier} gives the factor
\[
        \exp\{-c\operatorname{dist}(u_j-u_{\kappa(j)},2\pi\mathbb Z)^2\},
        \qquad u_0=0 .
\]
Multiplying over the disjoint blocks gives
\[
        |\phi(u)|\le
        \exp\left[-c\sum_{j\in G}A_j
        \operatorname{dist}(u_j-u_{\kappa(j)},2\pi\mathbb Z)^2\right].
\]
Because every directed edge goes from $j$ to a larger index, or to the
anchor $0$, each path reaches $0$ after at most $r$ steps.  Along such a
path the weights only increase.  Therefore, for $u_0=0$,
\[
        \sum_{j\in G}A_j u_j^2
        \le C_{r,\eta}\sum_{j\in G}A_j(u_j-u_{\kappa(j)})^2 .
\]
The torus version follows from the same path argument.  For each $j$, choose
integer multiples of $2\pi$ successively along the directed path $P(j)$ from
$j$ to the anchor $0$.  The triangle inequality on the circle gives
\[
        \operatorname{dist}(u_j,2\pi\mathbb Z)
        \le \sum_{\ell\in P(j)}
        \operatorname{dist}(u_\ell-u_{\kappa(\ell)},2\pi\mathbb Z).
\]
Cauchy's inequality, together with monotonicity of the weights along this path,
gives the displayed comparison with distances to $2\pi\mathbb Z$.
Fourier inversion then gives
\[
        \Pp(W_G=z_G)
        \le C\int_{[-\pi,\pi]^G}
        \exp\left[-c\sum_{j\in G}A_ju_j^2\right]\dd u_G
        \le C\prod_{j\in G}A_j^{-1/2}.
\]
Bounded tilts distort all atoms in the finite support by bounded factors, so
the two-atom lower bounds persist after changing constants.
\end{proof}

\begin{lemma}[Terminal interval edge blocks]\label{lem:terminal-interval-edge-blocks}
Fix \(r\ge1\), \(0<c_*<C_*<\infty\), and \(c_0>0\).  In the terminal
interval notation above, fix \(j\in\{1,\ldots,r\}\) and suppose
\[
        c_*\le s_j=a_jt_j\le C_*,
        \qquad
        m_j(t)=\sum_{b\notin A}\Pp(T_b\in I_j)\ge c_0a_j.
\]
Then there are constants \(0<\alpha<\beta<\infty\) and \(c>0\), depending
only on \(r,c_*,C_*\), and \(c_0\), and a set of at least \(ca_j\) rates
\(b\in[\alpha a_j,\beta a_j]\setminus A\) such that
\[
        \Pp(T_b\in I_j)\ge c,
        \qquad
        \Pp(T_b\le t_j)\ge c.
\]
Consequently, after projecting any chosen subset of the interval coordinates
larger than \(j\) either to their own coordinates or to the anchor \(0\), a
further subblock of size at least \(ca_j\) gives two projected atoms
\(e_j\) and \(e_{\kappa(j)}\), where \(\kappa(j)>j\) is a retained
coordinate or \(\kappa(j)=0\), each with probability at least \(c\).
The conclusion remains valid under every interval-coordinate exponential tilt
whose coordinates are bounded by a fixed constant, after changing \(c\).
\end{lemma}

\begin{proof}
Choose \(\alpha>0\) small and \(\beta<\infty\) large.  The contribution to
\(m_j(t)\) from \(b<\alpha a_j\) is at most \(\alpha a_j+1\).  Since
\(t_j\ge c_*/a_j\), the contribution from \(b>\beta a_j\) is at most
\[
        \sum_{b>\beta a_j}e^{-bt_j}
        \le C a_j e^{-c_*\beta}.
\]
After first decreasing \(\alpha\) and then increasing \(\beta\), at least a
fixed positive multiple of \(a_j\) of the total interval mass comes from
\([\alpha a_j,\beta a_j]\).  This block contains \(O(a_j)\) rates, so at
least \(ca_j\) of them satisfy \(\Pp(T_b\in I_j)\ge c\).  For every such
rate,
\[
        bt_j\ge\alpha c_*,
        \qquad
        \Pp(T_b\le t_j)=1-e^{-bt_j}\ge1-e^{-\alpha c_*}>0.
\]
The event \(\{T_b\le t_j\}\) is the disjoint union of the later categories
\(I_{j+1},\ldots,I_{r+1}\).  For each selected rate, one of these finitely
many categories has probability bounded below.  A second pigeonhole argument
produces a common later category on a subblock of size \(ca_j\).  Projection
maps that category either to a retained atom \(e_{\kappa(j)}\) or to the anchor
\(0\).  Removing the fixed set \(A\) changes the block size by only \(O_r(1)\).
A bounded exponential tilt changes every atom probability in the finite
categorical support by a bounded multiplicative factor, which proves the final
assertion.
\end{proof}

\begin{lemma}[Projected good--bad interval bound]\label{lem:projected-good-bad-interval-bound}
Fix $r\ge1$, $\eta>0$, and a compact scaled chamber
$K\subset(0,\infty)^r$.  In the terminal interval notation above, suppose
$s=(a_1t_1,\ldots,a_rt_r)\in K$ and $t_1>\cdots>t_r$.  Suppose the deterministic targets satisfy
\[
        c_0a_j\le y_j\le C_0a_j,
        \qquad 1\le j\le r .
\]
Then, uniformly for such $t$ and $y$,
\[
        \Pp(Y(t)=y)
        \le C\prod_{j=1}^r a_j^{-1/2}
        \exp\left[-c\sum_{j=1}^r
        \min\left\{\frac{(y_j-m_j(t))^2}{a_j},a_j\right\}\right].
\]
The constants depend only on $r,\eta,K,c_0$, and $C_0$.
\end{lemma}

\begin{proof}
Write $d_j=y_j-m_j(t)$.  Since $s\in K$, finite-geometric estimates give
$m_j(t)=O_K(a_j)$ and hence $|d_j|\le M_Ka_j$.  Fix a small
$\varepsilon>0$.  Let
\[
        G=\{j: |d_j|\le\varepsilon a_j\},
        \qquad B=\{1,\ldots,r\}\setminus G .
\]
For $j\in G$ the target lower bound gives
$m_j(t)\ge(c_0-\varepsilon)a_j$; choose
$\varepsilon<c_0/2$, so Lemma \ref{lem:terminal-interval-edge-blocks}
applies to every good coordinate.

Use one combined exponential tilt.  Put
\[
        \lambda_j=
        \begin{cases}
        \delta d_j/a_j, & j\in G,\\
        \delta\,\operatorname{sgn}(d_j), & j\in B,
        \end{cases}
\]
where $\delta>0$ is a sufficiently small fixed constant.  Let
$K_t(\lambda)=\log\Ee\exp\{\lambda\cdot(Y(t)-m(t))\}$.
For every bounded tilt $\xi$, the categorical support and the finite-geometric
bound $\Var_\xi Y_j(t)\le C a_j$ give
\[
        \Var_\xi(v\cdot Y(t))\le C\sum_{j=1}^r a_jv_j^2.
\]
Therefore
\[
        K_t(\lambda)\le C\sum_{j=1}^r a_j\lambda_j^2 .
\]
On the other hand,
\[
        \lambda\cdot d
        \ge \delta\sum_{j\in G}\frac{d_j^2}{a_j}
            +\delta\varepsilon\sum_{j\in B}a_j .
\]
Choosing $\delta$ small enough yields
\[
        -\lambda\cdot d+K_t(\lambda)
        \le -c\sum_{j\in G}\frac{d_j^2}{a_j}
             -c\sum_{j\in B}a_j .
\]
Under this bounded tilt all categorical probabilities are changed by bounded
multiplicative factors.  Project the tilted array onto the good coordinates
and ignore the bad-coordinate constraints.  For each $j\in G$, Lemma \ref{lem:terminal-interval-edge-blocks} gives a candidate block $B_j$ of at least $c_0a_j$
rates with two projected atoms $e_j$ and $e_{\kappa(j)}$, where
$\kappa(j)>j$ is another good coordinate or $\kappa(j)=0$.  We now extract
disjoint subblocks rigorously.  Process the indices of $G$ in increasing order.
Since $a_{i+1}/a_i\ge1+\eta$,
\[
        \sum_{i<j}a_i\le C_{r,\eta}a_j.
\]
Choose $c'>0$ so small that
$c'(1+C_{r,\eta})\le c_0$.  After subblocks of sizes $c'a_i$ have been chosen
for all earlier good indices $i<j$, at most
$c'\sum_{i<j}a_i\le c'C_{r,\eta}a_j$ elements of $B_j$ have been used.
At least $c'a_j$ unused elements of $B_j$ therefore remain.  Continuing
greedily produces disjoint subblocks of size at least $c'a_j$ for every
$j\in G$.
Lemma \ref{lem:projected-forward-edge-concentration} gives
\[
        \Pp_\lambda(Y_G(t)=y_G)\le C\prod_{j\in G}a_j^{-1/2}.
\]
Consequently
\begin{align*}
        \Pp(Y(t)=y)
        &\le \exp\{-\lambda\cdot d+K_t(\lambda)\}
        \Pp_\lambda(Y_G(t)=y_G)  \\
        &\le C\prod_{j\in G}a_j^{-1/2}
        \exp\left[-c\sum_{j\in G}\frac{d_j^2}{a_j}
                  -c\sum_{j\in B}a_j\right].
\end{align*}
For $j\in B$, the factor $e^{-ca_j}$ absorbs the missing local factor:
$e^{-ca_j}\le C a_j^{-1/2}e^{-c'a_j}$.  Since for $j\in G$ we have
$|d_j|\le\varepsilon a_j$, this is the desired bound, after changing
constants.

\end{proof}


\begin{lemma}[Actual-mean cumulative Chernoff bounds]\label{lem:actual-mean-lower-cumulative-chernoff}
Fix \(r\ge1\).  Let \(B\subset\{1,\ldots,r\}\), and put
\[
        \mu_j=\Ee Z_j(t),\qquad j\in B.
\]
There are constants \(\delta_0,C>0\), depending only on \(r\), such that,
for every \(|\theta|\le\delta_0\),
\[
\log \Ee\exp\left\{\theta\sum_{j\in B}(Z_j(t)-\mu_j)\right\}
        \le C\theta^2\sum_{j\in B}\mu_j .
\]
After decreasing \(\delta_0\), if necessary, fix
\(0<\delta\le\delta_0\) and \(C_0<\infty\).  There are
\(c=c(\delta,r,C_0)>0\) and \(L_0=L_0(\delta,r,C_0)\) such that, whenever
\(L\ge L_0\), both of the following assertions hold.

If \(z_j\le C_0L^{-1}\mu_j\) for every \(j\in B\), then
\[
        -\sum_{j\in B}(-\delta)(z_j-\mu_j)
        +\log \Ee\exp\left\{-\delta\sum_{j\in B}(Z_j(t)-\mu_j)\right\}
        \le -c\sum_{j\in B}\mu_j .
\]
If \(\mu_j\le C_0L^{-1}z_j\) for every \(j\in B\), then
\[
        -\sum_{j\in B}\delta(z_j-\mu_j)
        +\log \Ee\exp\left\{\delta\sum_{j\in B}(Z_j(t)-\mu_j)\right\}
        \le -c\sum_{j\in B}z_j .
\]
In particular, the first exponent is at most
\(-cL\sum_{j\in B}a_j\) if \(\mu_j\ge c_0La_j\), while the second is at most
\(-c\sum_{j\in B}a_j\) if \(z_j\ge c_0a_j\).
\end{lemma}

\begin{proof}
For a single rate \(b\notin A\), define
\[
        R_b=\sum_{j\in B}\one_{\{T_b>t_j\}},
        \qquad m_b=\Ee R_b .
\]
Then \(0\le R_b\le r\), and
\[
        \sum_{b\notin A}m_b=\sum_{j\in B}\mu_j .
\]
For \(|\theta|\le1\), let
\[
        K_b(\theta)=\log \Ee e^{\theta(R_b-m_b)} .
\]
Then \(K_b(0)=K_b'(0)=0\), and
\[
        K_b''(\theta)=\Var_\theta(R_b)
        \le \Ee_\theta R_b^2
        \le r\,\Ee_\theta R_b .
\]
Since \(0\le R_b\le r\), a bounded tilt changes the expectation of the
nonnegative variable \(R_b\) by at most a factor depending only on \(r\):
\[
        \Ee_\theta R_b
        =\frac{\Ee R_be^{\theta R_b}}{\Ee e^{\theta R_b}}
        \le e^{2r}m_b,
        \qquad |\theta|\le1 .
\]
Therefore \(K_b''(\theta)\le C_rm_b\) for \(|\theta|\le1\), and Taylor's
formula gives \(K_b(\theta)\le C_r\theta^2m_b\).  Summing over
\(b\notin A\) proves the first displayed inequality for both signs of
\(\theta\).

For the lower-tail assertion, use the preceding bound and
\(z_j\le C_0L^{-1}\mu_j\):
\[
\begin{aligned}
&-\sum_{j\in B}(-\delta)(z_j-\mu_j)
 +\log \Ee\exp\left\{-\delta\sum_{j\in B}(Z_j(t)-\mu_j)\right\} \\
&\qquad\le
        -\delta\left(1-\frac{C_0}{L}\right)\sum_{j\in B}\mu_j
        +C\delta^2\sum_{j\in B}\mu_j .
\end{aligned}
\]
Decrease \(\delta_0\), if necessary, so that \(C\delta_0\le1/4\), and then
choose \(L_0\) so that \(C_0/L\le1/4\).  The last display is at most
\(-\delta\sum_{j\in B}\mu_j/2\).

For the upper-tail assertion, the same cumulant bound and
\(\mu_j\le C_0L^{-1}z_j\) give
\[
\begin{aligned}
&-\sum_{j\in B}\delta(z_j-\mu_j)
 +\log \Ee\exp\left\{\delta\sum_{j\in B}(Z_j(t)-\mu_j)\right\} \\
&\qquad\le
        -\delta\sum_{j\in B}z_j
        +(\delta+C\delta^2)\sum_{j\in B}\mu_j\\
&\qquad\le
        -\delta\left(1-\frac{2C_0}{L}\right)
        \sum_{j\in B}z_j,
\end{aligned}
\]
after decreasing \(\delta_0\) once more.  Enlarging \(L_0\) makes this at
most \(-\delta\sum_{j\in B}z_j/2\).  The two final scale bounds follow
immediately from the corresponding assumptions.
\end{proof}

\begin{lemma}[Mixed cumulative--interval localization]
\label{lem:projected-cumulative-tail-localization}
Fix $r\ge1$ and $\eta>0$.  There is $L_0=L_0(r,\eta)$ such that the
following holds for every $L\ge L_0$.  In the terminal interval notation, let
$s_j=a_jt_j$ and suppose $s\notin[L^{-1},L]^r$.  Put
\[
        B_-(s)=\{j:s_j<L^{-1}\},\qquad
        B_+(s)=\{j:s_j>L\},\qquad
        G(s)=\{1,\ldots,r\}\setminus(B_-(s)\cup B_+(s)).
\]
On every region on which these three sets are fixed,
\[
\begin{aligned}
        \Pp(Y(t)=y)
        &\le C_L\prod_{j=1}^r a_j^{-1/2}
        \exp\left[-c\sum_{j\in B_-}|H_j(t)|
                   -c(\log L)\sum_{j\in B_+}a_j\right]\\
        &\quad\times
        \exp\left[-c_L\sum_{j\in G(s)}
        \min\left\{\frac{H_j(t)^2}{a_j},a_j\right\}\right].
\end{aligned}\tag{7.3}
\]
Here $c>0$ is independent of $L$, while $C_L<\infty$ and $c_L>0$
depend only on $L,r$, and $\eta$.  Consequently,
\[
\begin{aligned}
        \Pp(Y(t)=y)
        &\le C_L\prod_{j=1}^r a_j^{-1/2}
        \exp\left[-cL\sum_{j\in B_-}a_j
                   -c(\log L)\sum_{j\in B_+}a_j\right]\\
        &\quad\times
        \exp\left[-c_L\sum_{j\in G(s)}
        a_j\min\{(s_j-s_j^*)^2,1\}\right].
\end{aligned}\tag{7.4}
\]
\end{lemma}

\begin{proof}
Throughout the proof the offender sets are fixed.  Write $B_-=B_-(s)$,
$B_+=B_+(s)$, and $G=G(s)$.  Finite-geometric estimates give, uniformly in
the admissible terminal range,
\[
\begin{array}{lll}
 j\in B_-:& \mu_j(t_j)\ge cLa_j,& H_j(t)=z_j-\mu_j(t_j)<0,\\[2mm]
 j\in B_+:& \mu_j(t_j)\le Ca_j/L,& 0<H_j(t)\le z_j\le a_j,
\end{array}
\tag{7.5}
\]
for $L$ sufficiently large.  In particular,
$|H_j(t)|\ge cLa_j$ on $B_-$ and $H_j(t)\ge ca_j$ on $B_+$.
For $j\in G$, one has $s_j\in[L^{-1},L]$, and therefore
\[
        m_j(t)\le\mu_j(t_j)\le C_La_j,
        \qquad |d_j(t)|\le M_La_j.
        \tag{7.6}
\]

Let
\[
        K_t(\theta)=\log\Ee\exp\{\theta\cdot(Y(t)-m(t))\}
\]
be the centered interval-coordinate cumulant.  If
$\gamma=(\gamma_1,\ldots,\gamma_r)$ is a cumulative-coordinate vector, put
\[
        (\mathcal T\gamma)_k=\sum_{j\ge k}\gamma_j.
\]
Since $Z_j=Y_1+\cdots+Y_j$ and $d_j=H_j-H_{j-1}$,
\[
        (\mathcal T\gamma)\cdot d=\gamma\cdot H.
        \tag{7.7}
\]
Choose $\delta>0$ so small that $3\delta$ is admissible in Lemma
\ref{lem:actual-mean-lower-cumulative-chernoff} and
$3\delta r<1/2$.  Define
\[
        \gamma_j^-=-\delta\one_{\{j\in B_-\}},
        \qquad
        \gamma_j^+=\delta(\log L)\one_{\{j\in B_+\}}.
\]

We also introduce a local interval tilt on $G$.  Ratio separation implies that
there is $c_\eta>0$ such that $y_j\ge c_\eta a_j$ for all $j$ and all
sufficiently large $n$.  Fix $0<\varepsilon<c_\eta/4$ and put
\[
        G_0=\{j\in G:|d_j|\le\varepsilon a_j\},
        \qquad G_1=G\setminus G_0.
\]
For a constant $\delta_L>0$ to be chosen below, set
\[
        \lambda_j=
        \begin{cases}
        \delta_L d_j/a_j,&j\in G_0,\\
        \delta_L\operatorname{sgn}(d_j),&j\in G_1,\\
        0,&j\notin G.
        \end{cases}
\]
Finally, let
\[
        \theta=\mathcal T\gamma^-+\mathcal T\gamma^++\lambda.
\]
Generalized H\"older inequality gives the precise combined-cumulant estimate
\[
\begin{aligned}
        K_t(\theta)
        &\le\frac13K_t(3\mathcal T\gamma^-)
             +\frac13K_t(3\mathcal T\gamma^+)
             +\frac13K_t(3\lambda).
\end{aligned}\tag{7.8}
\]
Thus, by (7.7),
\[
\begin{aligned}
-\theta\cdot d+K_t(\theta)
&\le\frac13\{-3\gamma^-\cdot H+K_t(3\mathcal T\gamma^-)\}\\
&\quad+\frac13\{-3\gamma^+\cdot H+K_t(3\mathcal T\gamma^+)\}\\
&\quad+\frac13\{-3\lambda\cdot d+K_t(3\lambda)\}.
\end{aligned}\tag{7.9}
\]
We estimate the three braces separately.

For $B_-$, (7.5) implies $z_j\le C L^{-1}\mu_j(t_j)$.  Lemma
\ref{lem:actual-mean-lower-cumulative-chernoff}, applied with parameter
$3\delta$, gives
\[
        -3\gamma^-\cdot H+K_t(3\mathcal T\gamma^-)
        \le-c\sum_{j\in B_-}\mu_j(t_j)
        \le-c\sum_{j\in B_-}|H_j(t)|.
        \tag{7.10}
\]
The same bound also retains the weaker loss
$-cL\sum_{j\in B_-}a_j$.

For $B_+$, put
\[
        R_b^+=\sum_{j\in B_+}\one_{\{T_b>t_j\}},
        \qquad 0\le R_b^+\le r.
\]
Since the summands indexed by $b\notin A$ are independent,
\[
\begin{aligned}
K_t(3\mathcal T\gamma^+)
&=\sum_{b\notin A}
 \log\Ee\exp\{3\delta(\log L)(R_b^+-\Ee R_b^+)\}\\
&\le \sum_{b\notin A}
 \log\Ee\exp\{3\delta(\log L)R_b^+\}\\
&\le L^{3\delta r}\sum_{b\notin A}\Pp(R_b^+>0)\\
&\le L^{3\delta r}\sum_{j\in B_+}\mu_j(t_j)
 \le C L^{-1/2}\sum_{j\in B_+}a_j.
\end{aligned}
\tag{7.11}
\]
On the other hand, (7.5) gives
\[
        -3\gamma^+\cdot H
        \le-c(\log L)\sum_{j\in B_+}a_j.
\]
After increasing $L_0$, (7.11) is absorbed by this linear gain, and hence
\[
        -3\gamma^+\cdot H+K_t(3\mathcal T\gamma^+)
        \le-c(\log L)\sum_{j\in B_+}a_j.
        \tag{7.12}
\]

For the local tilt, let $v$ vanish outside $G$.  Under every bounded
interval-coordinate tilt $\xi$, a categorical summand contributes at most one
coordinate, and therefore
\[
        \Var_\xi(v\cdot Y(t))
        \le\sum_{j\in G}v_j^2\Ee_\xi Y_j(t)
        \le C_L\sum_{j\in G}a_jv_j^2,
        \tag{7.13}
\]
where the last inequality follows from (7.6) and bounded distortion of atom
probabilities.  Taylor's formula for the cumulant gives
\[
        K_t(3\lambda)\le C_L\sum_{j\in G}a_j\lambda_j^2.
\]
Choose $\delta_L>0$ sufficiently small.  Then
\[
        -3\lambda\cdot d+K_t(3\lambda)
        \le-c_L\sum_{j\in G}
        \min\left\{\frac{d_j^2}{a_j},a_j\right\}.
        \tag{7.14}
\]
Indeed, on $G_0$ the linear term is a fixed multiple of $d_j^2/a_j$,
whereas on $G_1$ it is at least a fixed multiple of $a_j$.

We next recover the product local factor under the single total tilt $\theta$.
For fixed $L$, its interval coordinates are bounded by a constant depending
only on $L$ and $r$.  If $j\in G_0$, then
\[
        m_j(t)\ge y_j-\varepsilon a_j\ge(c_\eta/2)a_j.
\]
Lemma \ref{lem:terminal-interval-edge-blocks}, with
$s_j\in[L^{-1},L]$, supplies a candidate forward-edge block of order $a_j$.
Process the indices of $G_0$ increasingly and choose small subblocks greedily.
Since $a_{i+1}/a_i\ge1+\eta$ and $r$ is fixed,
\[
        \sum_{i<j}a_i\le C_{r,\eta}a_j,
\]
so disjoint subblocks of size $c_{L,r,\eta}a_j$ can be retained for every
$j\in G_0$.  Boundedness of the total tilt preserves the two-atom lower
bounds.  Lemma \ref{lem:projected-forward-edge-concentration}, applied after
projection onto $G_0$, yields
\[
        \Pp_\theta(Y_{G_0}(t)=y_{G_0})
        \le C_L\prod_{j\in G_0}a_j^{-1/2}.
        \tag{7.15}
\]
The exact change-of-measure identity and (7.9)--(7.15) now give
\[
\begin{aligned}
\Pp(Y(t)=y)
&\le C_L\prod_{j\in G_0}a_j^{-1/2}
 \exp\left[-c\sum_{j\in B_-}|H_j|
            -c(\log L)\sum_{j\in B_+}a_j\right]\\
&\quad\times
 \exp\left[-c_L\sum_{j\in G}
 \min\left\{\frac{d_j^2}{a_j},a_j\right\}\right].
\end{aligned}\tag{7.16}
\]
For $j\in G_1$, the last exponent contains $-c_La_j$; for $j\in B_-$ and
$j\in B_+$, (7.5) and the preceding tail losses contain the same fixed-order
exponential cost.  These factors absorb all missing $a_j^{-1/2}$ terms, after
changing constants.

It remains to pass from interval residuals to cumulative residuals without
losing boundary terms.  The assumptions of Lemma
\ref{lem:boundary-truncated-chain} hold with $A_j=a_j$ and $M=M_L$: (7.6)
controls $d_j$ on $G$, while (7.5) gives $|H_j|\le a_j$ on $B_+$.  Split a
fixed fraction of each offender loss in (7.16) for this comparison.  Since
$\log L$ is large, Lemma \ref{lem:boundary-truncated-chain} yields
\[
\begin{aligned}
&c_L\sum_{j\in G}\min\left\{\frac{d_j^2}{a_j},a_j\right\}
 +c\sum_{j\in B_-}|H_j|
 +c(\log L)\sum_{j\in B_+}a_j\\
&\qquad\ge
 c_L'\sum_{j\in G}\min\left\{\frac{H_j^2}{a_j},a_j\right\}
 +c'\sum_{j\in B_-}|H_j|
 +c'(\log L)\sum_{j\in B_+}a_j.
\end{aligned}
\tag{7.17}
\]
Combining (7.16) and (7.17) proves (7.3).

Finally, on $G$ the scaled coordinate $s_j$ belongs to $[L^{-1},L]$.
Uniform finite-geometric derivative estimates and strict monotonicity give
\[
        |H_j(t)|=|\mu_j(t_j^*)-\mu_j(t_j)|
        \ge c_La_j\min\{|s_j-s_j^*|,1\}.
        \tag{7.18}
\]
Together with (7.5), this converts (7.3) into (7.4).
\end{proof}


\begin{lemma}[Anisotropic tilted local upper bound]\label{lem:anisotropic-tilted-upper}
Let \(W=\sum_{\nu\in\mathcal I}W_\nu\) be as in Lemma \ref{lem:coord-fourier}.  Assume that the hypotheses of Lemma \ref{lem:coord-fourier} hold uniformly under all tilts \(\|\lambda\|_\infty\le\Lambda\).  Assume also that, for those tilted laws,
\[
        \Var_\lambda(v\cdot W)
        \le C_0\sum_{j=1}^r A_jv_j^2,
        \qquad v\in\Rr^r .
\]
Then there are constants \(C,c>0\), depending only on the constants in the assumptions, such that for every lattice point \(z\), with
\(h=z-\Ee W\),
\[
        \Pp(W=z)
        \le C\prod_{j=1}^r A_j^{-1/2}
        \exp\left[-c\sum_{j=1}^r
        \min\left\{\frac{h_j^2}{A_j}, |h_j|\right\}\right].
\]
In particular, if \(|h_j|\le A_j\) for every \(j\), then the exponent may be written as
\(-c\sum_j h_j^2/A_j\).
\end{lemma}

\begin{proof}
Let
\[
        K(\lambda)=\log\Ee e^{\lambda\cdot(W-\Ee W)}.
\]
For \(\|\lambda\|_\infty\le\Lambda\), the covariance upper bound gives
\[
        K(\lambda)
        =\int_0^1(1-s)\lambda^T
        \Var_{s\lambda}(W)\lambda\dd s
        \le C\sum_{j=1}^r A_j\lambda_j^2 .
\]
Choose a sufficiently small \(\delta>0\), depending only on the constants in the assumptions and with \(\delta\le\Lambda\), and put
\[
        \lambda_j=\delta\,\mathrm{sgn}(h_j)
        \min\{|h_j|/A_j,1\}.
\]
By tilting and then applying Lemma \ref{lem:coord-fourier} under the tilted law,
\[
\begin{aligned}
        \Pp(W=z)
        &=\exp\{-\lambda\cdot h+K(\lambda)\}
          \Pp_\lambda(W=z) \\
        &\le C\prod_{j=1}^r A_j^{-1/2}
        \exp\{-\lambda\cdot h+K(\lambda)\}.
\end{aligned}
\]
For the chosen \(\lambda\),
\[
        \lambda\cdot h
        =\delta\sum_{j=1}^r |h_j|\min\{|h_j|/A_j,1\},
\]
while
\[
        K(\lambda)
        \le C\delta^2\sum_{j=1}^r
        A_j\min\{h_j^2/A_j^2,1\}.
\]
Taking \(\delta\) small enough makes the second display at most one half of the first, after changing constants.  This proves the claim.
\end{proof}

% END ARCHIVED BLOCK: shared-anisotropic-toolkit
% BEGIN ARCHIVED BLOCK: endpoint-regular-universality
\section{Endpoint-regular fixed-point universality}\label{sec:endpoint-fixed-universality}

We now promote the endpoint fixed-point prediction to a theorem.  The statement
is formulated in terminal rate coordinates.  Put
\[
        \lambda_{n,a}=\theta_{n,n+1-a},
        \qquad 1\le a\le n,
\]
so that terminal coordinate \(a\) is label \(n+1-a\).  Let
\[
        T_{n,a}\sim\Exp(\lambda_{n,a}),
        \qquad
        \tau_n(a)=1+\sum_{b\ne a}\one_{\{T_{n,b}>T_{n,a}\}}
        =n+1-\sigma_n(n+1-a).
\]
Thus $\tau_n$ is conjugate to $\sigma_n$ by the reversal permutation and has
the same cycle counts.  In Gumbel variables
\[
        Y_{n,a}=\log \lambda_{n,a}+G_a,
\]
the fixed-point indicator is
\[
        I_{n,a}=\one\left\{
        \sum_{b\ne a}\one_{\{Y_{n,b}<Y_{n,a}\}}=a-1
        \right\}.
\]
Indeed, label \(n+1-a\) is fixed exactly when there are \(a-1\) labels below
it in priority order.  We write
\[
        F_n=\sum_{a=1}^n I_{n,a}=C_1(\sigma_n)
\]
inside this section.

Multiplying all rates by a common deterministic factor adds the same constant to
all priorities and does not change the Luce order.  The following definition
therefore fixes a convenient normalization.

\begin{definition}[Uniform endpoint regular variation]\label{def:endpoint-rv}
The terminal rate array \((\lambda_{n,a})_{1\le a\le n}\) is
\emph{uniformly endpoint regularly varying with index \(\gamma>0\)} if there
exist positive constants \(s_n\), a constant \(C_0>0\), and a positive
measurable slowly varying function \(L\) at infinity such that the normalized
rates
\[
        \bar\lambda_{n,a}=s_n\lambda_{n,a}
\]
satisfy the following conditions.

\begin{enumerate}[label=(ER\arabic*)]
\item For every sequence \(A_n\to\infty\) with \(A_n=o(n)\),
\[
        \sup_{1\le a\le A_n}
        \left|
        \frac{\bar\lambda_{n,a}}
        {C_0(a/n)^\gamma L(n/a)}-1
        \right|\longrightarrow0.
\]
\item For every \(\delta>0\), there are constants
\(0<c_\delta<C_\delta<\infty\) such that
\[
        c_\delta\le \bar\lambda_{n,a}\le C_\delta,
        \qquad \delta n\le a\le n,
\]
for all sufficiently large \(n\).
\item The normalized terminal rates are nondecreasing:
\[
        \bar\lambda_{n,1}\le\bar\lambda_{n,2}
        \le\cdots\le\bar\lambda_{n,n}.
\]
\end{enumerate}
\end{definition}

A common multiplication of all rates does not change the Luce order.  We
therefore replace \(\lambda_{n,a}\) by \(\bar\lambda_{n,a}\) throughout this
section and suppress the bar.  We use repeatedly the uniform convergence
theorem and Potter bounds for slowly varying functions in the forms of
Bingham, Goldie, and Teugels \cite[Theorems~1.2.1 and~1.5.6]{BGT1987}.

For example, suppose \(\lambda_{n,a}=f(1-a/n)\), where
\[
        f(1-u)\sim C_0u^\gamma L(1/u),\qquad u\downarrow0,
\]
where \(f\) is bounded above and below away from zero on every compact
subinterval of \([0,1)\), and where the function
\(u\mapsto f(1-u)\) is nondecreasing on \((0,1]\).  Then the sampled array
satisfies Definition \ref{def:endpoint-rv}.  The global monotonicity in this
example is intentional: eventual endpoint monotonicity alone would not imply
condition \textup{(ER3)} for the full triangular array.

\begin{lemma}[The endpoint partial-sum bound is automatic]\label{lem:endpoint-partial-sum}
Under Definition \ref{def:endpoint-rv}, there is $c_\Sigma>0$ such that, for
all sufficiently large $n$,
\[
        \sum_{b=1}^a\lambda_{n,b}
        \ge c_\Sigma a\lambda_{n,a},
        \qquad 1\le a\le n.
\]
\end{lemma}

\begin{proof}
Suppose otherwise.  Then there are $n_k\to\infty$ and $a_k\in[n_k]$ such that
\[
        \frac{\sum_{b\le a_k}\lambda_{n_k,b}}
        {a_k\lambda_{n_k,a_k}}\longrightarrow0.
\]
Pass to a subsequence.  If $(a_k)$ is bounded, the ratio is at least $1/a_k$,
a contradiction.  Assume therefore that $a_k\to\infty$.

If $a_k/n_k\to0$, uniform endpoint regular variation and Potter bounds give
\[
        \frac1{a_k}\sum_{b=1}^{a_k}
        \frac{\lambda_{n_k,b}}{\lambda_{n_k,a_k}}
        \longrightarrow\int_0^1x^\gamma\dd x
        =\frac1{\gamma+1}.
\]
For completeness, on $\varepsilon a_k\le b\le a_k$ this is an ordinary
Riemann sum by the uniform convergence theorem for slowly varying functions.
On $b<\varepsilon a_k$, Potter's bound gives
$\lambda_{n_k,b}/\lambda_{n_k,a_k}\le C(b/a_k)^{\gamma/2}$, so the omitted
part is $O(\varepsilon^{1+\gamma/2})$.  Letting $\varepsilon\downarrow0$ proves
the display, again a contradiction.

The remaining case has $a_k\ge\delta n_k$ for some $\delta>0$.  By
monotonicity and (ER2),
\[
        \sum_{b=1}^{a_k}\lambda_{n_k,b}
        \ge\sum_{b=\lceil a_k/2\rceil}^{a_k}\lambda_{n_k,b}
        \ge \frac{a_k}{2}c_{\delta/2},
\]
whereas $\lambda_{n_k,a_k}\le C_\delta$.  The ratio is therefore bounded below
by $c_{\delta/2}/(2C_\delta)>0$, the final contradiction.
\end{proof}


Set
\[
        A_\gamma=\Gamma(1+1/\gamma),
        \qquad
        K_\gamma=A_\gamma^\gamma,
        \qquad
        c_\gamma=\gamma K_\gamma e^{-K_\gamma}.
\]

For fixed \(r\ge1\) and \(\eta>0\), define the admissible family
\[
\mathfrak E_N(r,\eta)=
\left\{
(n,a_1,\ldots,a_r):
N\le a_1<\cdots<a_r\le n/N,\quad
\frac{a_{j+1}}{a_j}\ge1+\eta
\right\}.
\]
Here and below, a statement asserted uniformly for
\(1\ll a_1<\cdots<a_r\ll n\) means that its supremum over
\(\mathfrak E_N(r,\eta)\) tends to zero as \(N\to\infty\).  This convention
makes all endpoint uniformity statements independent of a particular choice of
intermediate cutoff sequences.

\begin{lemma}[Weighted regularly varying terminal sums]
\label{lem:endpoint-rv-sums}
Assume Definition \ref{def:endpoint-rv}.  Let \(a=a_n\to\infty\) with
\(a=o(n)\).  For every fixed integer \(m\ge0\), put
\[
J_{\gamma,m}(\kappa)
 =\int_0^\infty x^{m\gamma}e^{-\kappa x^\gamma}\dd x
 =\frac1\gamma\Gamma\!\left(m+\frac1\gamma\right)
  \kappa^{-m-1/\gamma}.
\]
Then, uniformly for \(\kappa\) in compact subintervals of \((0,\infty)\),
\[
\sum_{b=1}^n
\left(\frac{\lambda_{n,b}}{\lambda_{n,a}}\right)^m
\exp\left(-\kappa\frac{\lambda_{n,b}}{\lambda_{n,a}}\right)
 =aJ_{\gamma,m}(\kappa)(1+o(1)).
\tag{8.1}\label{eq:endpoint-weighted-sums}
\]
The same assertion holds after deleting any fixed number of summands, uniformly
in the deleted indices.  In particular,
\[
\sum_{b=1}^n
\exp\left(-\kappa\frac{\lambda_{n,b}}{\lambda_{n,a}}\right)
 =aA_\gamma\kappa^{-1/\gamma}(1+o(1)),
\tag{8.2}\label{eq:endpoint-sum-zero}
\]
\[
\sum_{b=1}^n
\frac{\lambda_{n,b}}{\lambda_{n,a}}
\exp\left(-\kappa\frac{\lambda_{n,b}}{\lambda_{n,a}}\right)
 =\frac{aA_\gamma}{\gamma}
 \kappa^{-1-1/\gamma}(1+o(1)),
\tag{8.3}\label{eq:endpoint-sum-one}
\]
and, with
\(p_{n,b}(\kappa)=
\exp(-\kappa\lambda_{n,b}/\lambda_{n,a})\),
\[
\sum_{b=1}^n p_{n,b}(\kappa)(1-p_{n,b}(\kappa))
 =aA_\gamma\kappa^{-1/\gamma}
 (1-2^{-1/\gamma})(1+o(1)).
\tag{8.4}\label{eq:endpoint-variance-sum}
\]
\end{lemma}

\begin{proof}
Fix a compact interval
\([\kappa_-,\kappa_+]\Subset(0,\infty)\), and put
\[
        z_n=\frac na\longrightarrow\infty,
        \qquad
        R_n=\bigl(\log(e+z_n)\bigr)^{4/\gamma},
        \qquad
        B_n=\lceil aR_n\rceil.
\]
Since \(R_n=o(z_n)\), one has \(B_n=o(n)\).  Definition
\ref{def:endpoint-rv} therefore applies uniformly to all indices at most
\(B_n\).

Fix \(0<\varepsilon<1<R<\infty\).  On
\(\varepsilon a\le b\le Ra\), the uniform convergence theorem for slowly
varying functions gives
\[
        \frac{\lambda_{n,b}}{\lambda_{n,a}}
        =\left(\frac ba\right)^\gamma(1+o(1))
\]
uniformly.  Hence the contribution of this range to the left-hand side of
\eqref{eq:endpoint-weighted-sums}, divided by \(a\), converges uniformly in
\(\kappa\in[\kappa_-,\kappa_+]\) to
\[
        \int_\varepsilon^R x^{m\gamma}e^{-\kappa x^\gamma}\dd x.
\]
The range \(b<\varepsilon a\) contains at most \(\varepsilon a\) terms, and
\[
        \sup_{x\ge0,\,\kappa\in[\kappa_-,\kappa_+]}
        x^m e^{-\kappa x}<\infty.
\]
Its normalized contribution is therefore \(O_m(\varepsilon)\).

It remains to control the upper tail without applying regular variation at a
fixed positive fraction of \(n\).  Potter bounds, with exponent
\(\gamma/2\), imply that for all sufficiently large \(n\), uniformly for
\(Ra<b\le B_n\),
\[
 c\left(\frac ba\right)^{\gamma/2}
 \le \frac{\lambda_{n,b}}{\lambda_{n,a}}
 \le C\left(\frac ba\right)^{3\gamma/2}.
\]
Consequently the normalized contribution of this range is at most
\[
 C_m\int_R^\infty
 x^{3m\gamma/2}\exp(-c\kappa_-x^{\gamma/2})\dd x+o(1),
\]
which tends to zero as \(R\to\infty\).

For \(b>B_n\), monotonicity and the preceding Potter lower bound give
\[
 \frac{\lambda_{n,b}}{\lambda_{n,a}}
 \ge \frac{\lambda_{n,B_n}}{\lambda_{n,a}}
 \ge cR_n^{\gamma/2}
 =c\bigl(\log(e+z_n)\bigr)^2.
\]
Therefore
\[
\begin{aligned}
&\frac1a\sum_{b>B_n}
 \left(\frac{\lambda_{n,b}}{\lambda_{n,a}}\right)^m
 e^{-\kappa\lambda_{n,b}/\lambda_{n,a}}\\
&\qquad\le z_n
 \sup_{x\ge c(\log(e+z_n))^2}x^m e^{-\kappa_-x}
 \longrightarrow0.
\end{aligned}
\]
Letting first \(n\to\infty\), then \(\varepsilon\downarrow0\) and
\(R\uparrow\infty\), proves \eqref{eq:endpoint-weighted-sums}.  A fixed
deleted summand is bounded uniformly by
\(\sup_{x\ge0}x^m e^{-\kappa_-x}\), so finitely many deletions change the sum
by \(O(1)=o(a)\).  The special cases
\eqref{eq:endpoint-sum-zero}--\eqref{eq:endpoint-sum-one} follow from
\(J_{\gamma,0}(\kappa)=A_\gamma\kappa^{-1/\gamma}\) and
\(J_{\gamma,1}(\kappa)=A_\gamma\kappa^{-1-1/\gamma}/\gamma\).
Finally, \eqref{eq:endpoint-variance-sum} is the difference of
\eqref{eq:endpoint-sum-zero} at \(\kappa\) and at \(2\kappa\).
\end{proof}

\begin{lemma}[Endpoint product envelope]\label{lem:endpoint-rv-product-bound}
Assume Definition \ref{def:endpoint-rv}.  For every finite set
\(A\subset[n]\),
\[
        \Pp(I_{n,a}=1\text{ for all }a\in A)
        \le C^{|A|}\prod_{a\in A}\frac1a,
\]
where $C$ depends only on the constant $c_\Sigma$ from Lemma \ref{lem:endpoint-partial-sum}.
\end{lemma}

\begin{proof}
Reveal the Luce order as a death process.  When the alive set is \(S\), the next
label deleted from \(S\) is \(b\in S\) with probability
\(\lambda_{n,b}/\sum_{c\in S}\lambda_{n,c}\).  The event \(I_{n,a}=1\) says that
terminal coordinate \(a\) is deleted at the transition at which the number of
alive labels drops from \(a\) to \(a-1\).  Condition on the history before that
transition.  If coordinate \(a\) is already dead, the conditional probability is
zero.  If it is alive, the alive set has size \(a\) and contains \(a\).  Since
the rates are nondecreasing, the total rate of any alive set of size \(a\) is at
least $\sum_{b=1}^a\lambda_{n,b}$, and Lemma \ref{lem:endpoint-partial-sum} gives
\[
        \frac{\lambda_{n,a}}{\sum_{c\in S}\lambda_{n,c}}
        \le \frac{1}{c_\Sigma a}.
\]
Order the elements of \(A\) decreasingly and condition successively before the
corresponding death levels.  Multiplying the conditional bounds proves the
claim.
\end{proof}


\subsection{Separated endpoint fixed-point asymptotics}
\label{subsec:endpoint-separated-fixed}

Fix \(r\ge1\) and \(\eta>0\), and let
\((n,a_1,\ldots,a_r)\in\mathfrak E_N(r,\eta)\).  Put
\[
        B=\{a_1,\ldots,a_r\},
        \qquad z_j=a_j-j,
        \qquad 1\le j\le r.
\]
For \(g=(g_1,\ldots,g_r)\), set
\[
        y_j(g)=\log\lambda_{n,a_j}+g_j,
        \qquad
        Z_j(g)=\sum_{b\notin B}
        \one_{\{Y_{n,b}<y_j(g)\}},
\]
and write
\[
        M_j(g_j)=\Ee Z_j(g)
        =\sum_{b\notin B}
        \exp\left(-e^{-g_j}
        \frac{\lambda_{n,b}}{\lambda_{n,a_j}}\right).
\]
The relevant chamber is
\[
        \mathscr C_n=\{g:y_1(g)<\cdots<y_r(g)\}.
\]
When \(g\in\mathscr C_n\), put \(y_0=-\infty\),
\(y_{r+1}=\infty\), and define interval counts
\[
        W_1=Z_1,
        \qquad W_j=Z_j-Z_{j-1}\quad(2\le j\le r).
\]
Thus a nonspecified priority contributes \(e_j\) to \(W\) when it lies in
\([y_{j-1},y_j)\), and contributes \(0\) when it lies in
\([y_r,\infty)\).  The target interval vector is
\[
        w_1=a_1-1,
        \qquad
        w_j=a_j-a_{j-1}-1\quad(2\le j\le r),
\]
and ratio separation gives
\[
        c_\eta a_j\le w_j\le a_j
        \qquad(1\le j\le r)
\tag{8.5}\label{eq:endpoint-interval-target-scale}
\]
for all sufficiently large \(N\).

\begin{lemma}[Endpoint Gumbel saddle geometry]
\label{lem:endpoint-gumbel-saddle}
For each \(j\), the equation
\[
        M_j(g_j^*)=z_j
\]
has a unique solution.  Uniformly over
\(\mathfrak E_N(r,\eta)\), as \(N\to\infty\),
\[
        g_j^*\longrightarrow g_\gamma:=-\log K_\gamma,
        \qquad
        M_j'(g_j^*)\sim\frac{a_j}{\gamma}.
\tag{8.6}\label{eq:endpoint-saddle-asymptotics}
\]
Moreover, for every fixed \(L<\infty\),
\[
        \sup_{|g|\le L}|M_j''(g)|\le C_La_j,
        \qquad
        c_La_j\le M_j'(g)\le C_La_j
\tag{8.7}\label{eq:endpoint-mean-derivative-bounds}
\]
uniformly whenever \(|g-g_j^*|\le1\).  There are constants
\(d_\eta>0\), \(\delta_\eta>0\), and \(N_0\) such that, for
\(N\ge N_0\),
\[
        y_{j+1}(g_{j+1})-y_j(g_j)\ge d_\eta
\tag{8.8}\label{eq:endpoint-threshold-gap}
\]
for every \(g\) satisfying
\(|g_j-g_j^*|\le\delta_\eta\), \(1\le j\le r\).  In particular, this
entire product neighborhood lies in \(\mathscr C_n\).
\end{lemma}

\begin{proof}
The function \(M_j\) is continuous and strictly increasing from \(0\) to
\(n-r\), so the saddle exists uniquely.  With \(\kappa=e^{-g}\), Lemma
\ref{lem:endpoint-rv-sums}, including its finite-deletion assertion, gives
uniformly on compact \(g\)-sets
\[
        M_j(g)=a_jA_\gamma e^{g/\gamma}(1+o(1)).
\tag{8.9}\label{eq:endpoint-mean-asymptotic}
\]
Differentiating term by term and applying
\eqref{eq:endpoint-weighted-sums} with \(m=1\) gives
\[
\begin{aligned}
        M_j'(g)
        &=\sum_{b\notin B}e^{-g}
        \frac{\lambda_{n,b}}{\lambda_{n,a_j}}
        \exp\left(-e^{-g}
        \frac{\lambda_{n,b}}{\lambda_{n,a_j}}\right)\\
        &=\frac{a_jA_\gamma}{\gamma}e^{g/\gamma}(1+o(1)).
\end{aligned}
\tag{8.10}\label{eq:endpoint-mean-prime-asymptotic}
\]
The saddle equation and \(z_j/a_j\to1\) imply
\(A_\gamma e^{g_j^*/\gamma}\to1\), which is equivalent to
\(g_j^*\to-\gamma\log A_\gamma=-\log K_\gamma\).  Substitution in
\eqref{eq:endpoint-mean-prime-asymptotic} gives the derivative asymptotic.
Since
\[
        M_j''(g)=\sum_{b\notin B}
        \left[
        \left(e^{-g}\frac{\lambda_{n,b}}{\lambda_{n,a_j}}\right)^2
        -e^{-g}\frac{\lambda_{n,b}}{\lambda_{n,a_j}}
        \right]
        e^{-e^{-g}\lambda_{n,b}/\lambda_{n,a_j}},
\]
\eqref{eq:endpoint-weighted-sums} with \(m=1,2\) proves the second-derivative
bound.  The lower and upper derivative bounds then follow from
\eqref{eq:endpoint-mean-prime-asymptotic} on a compact neighborhood of the
saddles.

It remains to separate the thresholds.  Uniform endpoint regular variation and
Potter bounds imply that there is \(c_\eta>0\) such that
\[
        \log\frac{\lambda_{n,a_{j+1}}}{\lambda_{n,a_j}}
        \ge2c_\eta
\tag{8.11}\label{eq:endpoint-rate-gap}
\]
for all sufficiently large \(N\).  To verify this uniformly, first restrict
\(a_{j+1}/a_j\) to a fixed compact interval
\([1+\eta,R]\) and use the uniform convergence theorem; for ratios larger
than \(R\), use the Potter lower bound
\(\lambda_{n,a_{j+1}}/\lambda_{n,a_j}
\ge c(a_{j+1}/a_j)^{\gamma/2}\), and then choose \(R\) large.
Because all \(g_j^*\) converge to the same constant, the saddle threshold gaps
are at least \(c_\eta\) for large \(N\).  Taking
\(\delta_\eta<c_\eta/4\) proves \eqref{eq:endpoint-threshold-gap}.
\end{proof}

\begin{lemma}[Endpoint interval edge blocks]
\label{lem:endpoint-interval-edge-blocks}
Fix \(L_0>0\), and consider any vector \(g\in\mathscr C_n\).  Write
\[
        \mathcal I_j(g)=[y_{j-1}(g),y_j(g))\quad(1\le j\le r),
        \qquad
        \mathcal I_0(g)=[y_r(g),\infty).
\]
Let
\[
        m_j(g)=\Ee W_j(g).
\]
For every \(c_0>0\), there are constants
\(0<\alpha<\beta<\infty\) and \(c>0\), depending only on
\(r,\eta,L_0,c_0\), with the following property.  If
\(|g_j|\le L_0\) and \(m_j(g)\ge c_0a_j\), then at least \(ca_j\) indices
\(b\in[\alpha a_j,\beta a_j]\setminus B\) satisfy
\[
        \Pp(Y_{n,b}\in\mathcal I_j(g))\ge c
\]
and, for one fixed
\(\kappa(j)\in\{j+1,\ldots,r,0\}\),
\[
        \Pp(Y_{n,b}\in\mathcal I_{\kappa(j)}(g))\ge c.
\]
Here \(\kappa(j)=0\) denotes the anchor interval.
The selected subblocks may be chosen disjoint as \(j\) varies, after reducing
\(c\).

In the saddle neighborhood from Lemma
\ref{lem:endpoint-gumbel-saddle}, one may instead choose disjoint blocks of
size at least \(ca_j\) on which the two adjacent atoms \(e_j\) and
\(e_{j+1}\) each have probability at least \(c\), with
\(e_{r+1}=0\).  All these conclusions remain valid, with modified constants,
under every interval-coordinate exponential tilt whose coordinates are bounded
by a fixed constant.
\end{lemma}

\begin{proof}
The contribution to \(m_j(g)\) from \(b<\alpha a_j\) is at most
\(\alpha a_j\).  Since membership in the \(j\)-th interval implies
\(Y_{n,b}<y_j\), the contribution from \(b>\beta a_j\) is bounded by
\[
        \sum_{b>\beta a_j}
        \exp\left(-e^{-g_j}
        \frac{\lambda_{n,b}}{\lambda_{n,a_j}}\right).
\]
The moving-cutoff proof of Lemma \ref{lem:endpoint-rv-sums}, with a Potter
bound on the intermediate tail and monotonicity beyond the moving cutoff,
shows that this is \(o_\beta(a_j)\), uniformly whenever \(|g_j|\le L_0\).  Choose \(\alpha\)
small and \(\beta\) large.  A positive multiple of \(a_j\) of the expected
interval mass then comes from \(b\in[\alpha a_j,\beta a_j]\).  This range
contains \(O(a_j)\) indices, so at least \(ca_j\) of them have
\(j\)-th interval probability at least \(c\).

On the same middle range, regular variation on fixed multiplicative intervals
keeps \(\lambda_{n,b}/\lambda_{n,a_j}\) between two positive constants.
Since \(g_j\) is bounded, the probability
\(\Pp(Y_{n,b}\ge y_j)\) is also bounded below.  Split this upper event among
the finitely many later intervals and the anchor; the pigeonhole principle
produces a common \(\kappa(j)>j\), or the anchor \(0\), on a subblock of order
\(a_j\).  Disjoint subblocks are extracted greedily.  Indeed,
\(a_{i+1}/a_i\ge1+\eta\) and \(r\) is fixed, so
\(\sum_{i<j}a_i\le C_{r,\eta}a_j\).

In the saddle neighborhood, choose a smaller multiplicative block
\([(1-\delta)a_j,(1+\delta)a_j]\), with \(\delta\) small enough that these
blocks are disjoint.  On this block the centered threshold
\(y_j-\log\lambda_{n,b}\) remains in a compact interval, and
\eqref{eq:endpoint-threshold-gap} gives uniformly positive intervals on both
sides of \(y_j\).  The standard Gumbel density is positive and continuous, so
the immediate lower and upper interval probabilities are both bounded below.
This gives the adjacent atoms.  A bounded exponential tilt multiplies all
probabilities in the finite categorical support by factors bounded above and
below, so every two-atom lower bound persists.
\end{proof}

\begin{lemma}[Endpoint Gumbel anisotropic local CLT]
\label{lem:endpoint-gumbel-llt}
Let \(\mathcal U_n=\{g:|g_j-g_j^*|\le\delta_\eta,
1\le j\le r\}\), where \(\delta_\eta\) is from Lemma
\ref{lem:endpoint-gumbel-saddle}.  Let
\[
        \Sigma(g)=\Var Z(g),
        \qquad
        D_a=\diag(a_1,\ldots,a_r).
\]
There are constants \(0<c<C<\infty\) such that, uniformly over
\(\mathfrak E_N(r,\eta)\) and \(g\in\mathcal U_n\),
\[
        cD_a\le\Sigma(g)\le CD_a
\tag{8.12}\label{eq:endpoint-covariance-comparison}
\]
in quadratic-form order.  For every fixed \(R<\infty\), uniformly for lattice
points \(q\in\mathbb Z^r\) satisfying
\[
        \|D_a^{-1/2}(q-M(g))\|\le R,
\]
\[
\begin{aligned}
\Pp(Z(g)=q)
&=(2\pi)^{-r/2}(\det\Sigma(g))^{-1/2}\\
&\quad\times
\exp\left[-\frac12(q-M(g))^T\Sigma(g)^{-1}(q-M(g))\right]
+o\left(\prod_{j=1}^r a_j^{-1/2}\right).
\end{aligned}
\tag{8.13}\label{eq:endpoint-gumbel-llt}
\]
The error is uniform over the displayed family.  Moreover, for every fixed
\(\nu<1/6\),
\[
\Pp(Z(g)=q)
\le C\prod_{j=1}^r a_j^{-1/2}
\exp\left[-c\sum_{j=1}^r
\frac{(q_j-M_j(g_j))^2}{a_j}\right]
\tag{8.14}\label{eq:endpoint-gumbel-local-upper}
\]
uniformly whenever
\(|q_j-M_j(g_j)|\le a_j^{1/2+\nu}\) for every \(j\).
\end{lemma}

\begin{proof}
For \(b\notin B\), let
\[
        V_b(g)=
        (\one_{\{Y_{n,b}<y_1\}},\ldots,
         \one_{\{Y_{n,b}<y_r\}}).
\]
These independent vectors have the cumulative categorical support
\[
        (1,\ldots,1),(0,1,\ldots,1),\ldots,(0,\ldots,0,1),0.
\]
The adjacent blocks in Lemma \ref{lem:endpoint-interval-edge-blocks} imply
that, for every \(u\in\mathbb R^r\), the \(j\)-th block contributes at least
\(ca_ju_j^2\) to \(\Var(u\cdot Z(g))\).  Hence the covariance lower bound in
\eqref{eq:endpoint-covariance-comparison} holds.  The upper bound follows from
\(\Var Z_j(g)=O(a_j)\), which is
\eqref{eq:endpoint-variance-sum}, and Cauchy's inequality.

The same adjacent blocks give, uniformly under bounded tilts,
\[
        |\Ee e^{iu\cdot(Z(g)-M(g))}|
        \le\exp\left[-c\sum_{j=1}^r
        a_j\operatorname{dist}(u_j,2\pi\mathbb Z)^2\right].
\tag{8.15}\label{eq:endpoint-fourier-decay}
\]
Indeed, the two adjacent cumulative atoms differ by the coordinate vector
\(e_j\), and the blocks are disjoint.  Thus the lattice generated by the
support differences is the full lattice \(\mathbb Z^r\).

Fourier inversion now gives the local expansion.  On the anisotropic small box
\(|u_j|\le Ra_j^{-1/2}\), boundedness of the categorical increments and the
covariance upper bound give
\[
\log\Ee e^{iu\cdot(Z-M)}
=-\frac12u^T\Sigma(g)u
+O\left((\max_j|u_j|)\sum_{j=1}^ra_ju_j^2\right).
\]
After the change of variables \(u=D_a^{-1/2}v\), the error is
\(O_R(a_1^{-1/2})=o(1)\).  The complement of a fixed \(v\)-ball is uniformly
negligible on the product scale by \eqref{eq:endpoint-fourier-decay}.  This
proves \eqref{eq:endpoint-gumbel-llt}.

For \eqref{eq:endpoint-gumbel-local-upper}, apply Lemma
\ref{lem:anisotropic-tilted-upper} with \(W=Z(g)\) and \(A_j=a_j\).  The
coordinatewise Fourier hypothesis under bounded tilts follows from the
adjacent blocks, and the required tilted covariance upper bound follows from
bounded distortion of the categorical probabilities and the argument proving
\eqref{eq:endpoint-covariance-comparison}.  Since
\(|q_j-M_j(g_j)|\le a_j^{1/2+\nu}=o(a_j)\), the truncated exponent in that
lemma is exactly \((q_j-M_j(g_j))^2/a_j\).  This proves
\eqref{eq:endpoint-gumbel-local-upper}.
\end{proof}

\begin{lemma}[Endpoint mixed cumulative--interval localization]
\label{lem:endpoint-mixed-localization}
There is \(L_0=L_0(r,\eta,\gamma)\) such that the following holds for
\(L\ge L_0\).  For \(g\in\mathscr C_n\), put
\[
\begin{aligned}
        B_R(g)&=\{j:g_j>L\},\\
        B_L(g)&=\{j:g_j<-L\},\\
        G(g)&=\{1,\ldots,r\}\setminus(B_R(g)\cup B_L(g)),
\end{aligned}
\]
and let
\[
        H_j=z_j-M_j(g_j),\qquad H_0=0,
        \qquad d_j=H_j-H_{j-1}.
\]
On every region on which the three index sets are fixed,
\[
\begin{aligned}
\Pp(Z(g)=z)
&\le C_L\prod_{j=1}^ra_j^{-1/2}
\exp\left[-c\sum_{j\in B_R}|H_j|
           -c\sum_{j\in B_L}a_j\right]\\
&\quad\times
\exp\left[-c_L\sum_{j\in G}
\min\left\{\frac{H_j^2}{a_j},a_j\right\}\right].
\end{aligned}
\tag{8.16}\label{eq:endpoint-mixed-cumulative-bound}
\]
Consequently,
\[
\begin{aligned}
\Pp(Z(g)=z)
&\le C_L\prod_{j=1}^ra_j^{-1/2}
\exp\left[-ce^{L/(2\gamma)}\sum_{j\in B_R}a_j
           -c\sum_{j\in B_L}a_j\right]\\
&\quad\times
\exp\left[-c_L\sum_{j\in G}
 a_j\min\{(g_j-g_j^*)^2,1\}\right].
\end{aligned}
\tag{8.17}\label{eq:endpoint-mixed-scaled-bound}
\]
Here \(c>0\) is independent of \(L\), while \(C_L,c_L\) may depend on
\(L,r,\eta,\gamma\).
\end{lemma}

\begin{proof}
Fix the offender sets and abbreviate them by \(B_R,B_L,G\).  From
\eqref{eq:endpoint-mean-asymptotic}, uniformly for large \(L\),
\[
\begin{array}{lll}
 j\in B_R:&M_j(g_j)\ge ce^{L/\gamma}a_j,&H_j<0,\\[1mm]
 j\in B_L:&M_j(g_j)\le Ce^{-L/\gamma}a_j,&0<H_j\le a_j.
\end{array}
\tag{8.18}\label{eq:endpoint-offender-geometry}
\]
For \(j\in G\), \(g_j\in[-L,L]\), so
\(m_j(g)=O_L(a_j)\) and \(|d_j|\le M_La_j\).

Let
\[
        K_g(\theta)=\log\Ee
        \exp\{\theta\cdot(W(g)-m(g))\}
\]
be the centered interval-coordinate cumulant.  For a cumulative-coordinate
vector \(\alpha\), define
\[
        (\mathcal T\alpha)_k=\sum_{j\ge k}\alpha_j.
\]
Because \(Z_j=W_1+\cdots+W_j\),
\[
        (\mathcal T\alpha)\cdot d=\alpha\cdot H.
\tag{8.19}\label{eq:endpoint-cumulative-interval-duality}
\]
Choose \(\delta>0\) so small that \(3\delta\) is admissible in Lemma
\ref{lem:actual-mean-lower-cumulative-chernoff}.  Set
\[
        \alpha_j^-=-\delta\one_{\{j\in B_R\}},
        \qquad
        \alpha_j^+=\delta\one_{\{j\in B_L\}}.
\]
For the nonoffending interval coordinates, fix
\(0<\varepsilon<c_\eta/4\), where \(c_\eta\) is from
\eqref{eq:endpoint-interval-target-scale}, and put
\[
        G_0=\{j\in G:|d_j|\le\varepsilon a_j\},
        \qquad G_1=G\setminus G_0.
\]
For a sufficiently small constant \(\delta_L>0\), define
\[
        \lambda_j=
        \begin{cases}
        \delta_Ld_j/a_j,&j\in G_0,\\
        \delta_L\operatorname{sgn}(d_j),&j\in G_1,\\
        0,&j\notin G.
        \end{cases}
\]
Use the single total interval-coordinate tilt
\[
        \Theta=\mathcal T\alpha^-+\mathcal T\alpha^++\lambda.
\]
Generalized H\"older inequality gives
\[
\begin{aligned}
K_g(\Theta)
&\le\frac13K_g(3\mathcal T\alpha^-)
 +\frac13K_g(3\mathcal T\alpha^+)
 +\frac13K_g(3\lambda),
\end{aligned}
\tag{8.20}\label{eq:endpoint-holder-cumulant}
\]
and therefore, by \eqref{eq:endpoint-cumulative-interval-duality},
\[
\begin{aligned}
-\Theta\cdot d+K_g(\Theta)
&\le\frac13\{-3\alpha^-\cdot H+K_g(3\mathcal T\alpha^-)\}\\
&\quad+\frac13\{-3\alpha^+\cdot H+K_g(3\mathcal T\alpha^+)\}\\
&\quad+\frac13\{-3\lambda\cdot d+K_g(3\lambda)\}.
\end{aligned}
\tag{8.21}\label{eq:endpoint-holder-exponent}
\]

The proof of Lemma
\ref{lem:actual-mean-lower-cumulative-chernoff} uses only independence,
nested indicator counts, and the bound that one summand contributes at most
\(r\) to the cumulative linear form.  It therefore applies verbatim to the
present lower-priority counts.  For \(B_R\),
\eqref{eq:endpoint-offender-geometry} gives
\(z_j\le Ce^{-L/\gamma}M_j(g_j)\).  The lower-tail part of Lemma
\ref{lem:actual-mean-lower-cumulative-chernoff}, applied with parameter
\(3\delta\), yields
\[
        -3\alpha^-\cdot H+K_g(3\mathcal T\alpha^-)
        \le-c\sum_{j\in B_R}M_j(g_j)
        \le-c\sum_{j\in B_R}|H_j|.
\tag{8.22}\label{eq:endpoint-right-offender-cost}
\]
For \(B_L\), one has
\(M_j(g_j)\le Ce^{-L/\gamma}z_j\).  The upper-tail part of the same lemma
gives
\[
        -3\alpha^+\cdot H+K_g(3\mathcal T\alpha^+)
        \le-c\sum_{j\in B_L}z_j
        \le-c\sum_{j\in B_L}a_j.
\tag{8.23}\label{eq:endpoint-left-offender-cost}
\]

Under every bounded interval tilt \(\xi\), the categorical support implies
\[
        \Var_\xi(v\cdot W(g))
        \le C_L\sum_{j\in G}a_jv_j^2
\]
for vectors supported on \(G\).  Taylor's formula for the cumulant and a
sufficiently small choice of \(\delta_L\) therefore give
\[
        -3\lambda\cdot d+K_g(3\lambda)
        \le-c_L\sum_{j\in G}
        \min\left\{\frac{d_j^2}{a_j},a_j\right\}.
\tag{8.24}\label{eq:endpoint-local-tilt-cost}
\]

All coordinates of the total tilt \(\Theta\) are bounded.  For
\(j\in G_0\),
\(m_j(g)\ge w_j-\varepsilon a_j\ge(c_\eta/2)a_j\).
Lemma \ref{lem:endpoint-interval-edge-blocks} supplies forward-edge blocks of
order \(a_j\), which may be chosen disjoint.  After projection onto
\(G_0\), a later atom whose coordinate is not retained becomes the anchor
atom.  Lemma \ref{lem:projected-forward-edge-concentration} therefore gives
under the single tilted law
\[
        \Pp_\Theta(W_{G_0}=w_{G_0})
        \le C_L\prod_{j\in G_0}a_j^{-1/2}.
\tag{8.25}\label{eq:endpoint-tilted-local-factor}
\]
The exact change-of-measure identity, together with
\eqref{eq:endpoint-holder-exponent}--\eqref{eq:endpoint-tilted-local-factor},
now yields the same bound as \eqref{eq:endpoint-mixed-cumulative-bound}, but
with interval residuals \(d_j\) in place of cumulative residuals on \(G\),
and initially only the local factors for \(G_0\).  Every index in \(G_1\),
\(B_R\), or \(B_L\) carries a fixed exponential loss of order \(a_j\), which
absorbs its missing factor \(a_j^{-1/2}\).

It remains to compare interval and cumulative residuals without discarding the
boundary values.  Apply Lemma \ref{lem:boundary-truncated-chain} with
\(A_j=a_j\), with its set \(B_-\) equal to the present \(B_R\), and with its
set \(B_+\) equal to the present \(B_L\).  Its assumptions hold by
\(|d_j|\le M_La_j\) on \(G\) and \(|H_j|\le a_j\) on \(B_L\).  Reserving a
fixed fraction of the offender costs in
\eqref{eq:endpoint-right-offender-cost}--\eqref{eq:endpoint-left-offender-cost}
therefore gives \eqref{eq:endpoint-mixed-cumulative-bound}.

Finally, on \(G\), the derivative bounds in
\eqref{eq:endpoint-mean-derivative-bounds}, monotonicity, and compactness imply
\[
        |H_j|\ge c_La_j\min\{|g_j-g_j^*|,1\}.
\tag{8.26}\label{eq:endpoint-mean-displacement}
\]
On \(B_R\), \eqref{eq:endpoint-offender-geometry} gives
\(|H_j|\ge ce^{L/\gamma}a_j\), which may be weakened to
\(ce^{L/(2\gamma)}a_j\).  These observations convert
\eqref{eq:endpoint-mixed-cumulative-bound} into
\eqref{eq:endpoint-mixed-scaled-bound}.
\end{proof}

\begin{lemma}[Endpoint Gumbel product localization]
\label{lem:endpoint-gumbel-localization}
Uniformly over \(\mathfrak E_N(r,\eta)\), the following assertions hold.

\begin{enumerate}[label=(\roman*)]
\item For every fixed compact cube \(K\subset\mathbb R^r\) containing all
saddle vectors for large \(N\), there are \(C_K,c_K>0\) such that, throughout
\(K\cap\mathscr C_n\),
\[
        \Pp(Z(g)=z)
        \le C_K\prod_{j=1}^ra_j^{-1/2}
        \exp\left[-c_K\sum_{j=1}^r
        a_j\min\{(g_j-g_j^*)^2,1\}\right].
\tag{8.27}\label{eq:endpoint-compact-localization}
\]
\item With \(\varphi(g)=e^{-g}e^{-e^{-g}}\),
\[
\lim_{L\to\infty}\limsup_{N\to\infty}
\sup_{\mathfrak E_N(r,\eta)}
\frac{
\displaystyle
\int_{\mathscr C_n\cap\{g\notin[-L,L]^r\}}
\prod_{j=1}^r\varphi(g_j)\Pp(Z(g)=z)\dd g
}{\prod_{j=1}^ra_j^{-1}}=0.
\tag{8.28}\label{eq:endpoint-noncompact-localization}
\]
\end{enumerate}
\end{lemma}

\begin{proof}
Choose \(L\) so that \(K\subset[-L,L]^r\).  With no offenders, the proof of
Lemma \ref{lem:endpoint-mixed-localization} gives
\eqref{eq:endpoint-compact-localization}; equivalently, apply
\eqref{eq:endpoint-mixed-cumulative-bound} with \(G=\{1,\ldots,r\}\).

For the tail, partition the integration region according to the finitely many
sets \(B_R,B_L,G\), and use \eqref{eq:endpoint-mixed-scaled-bound}.  Every
nonoffending coordinate contributes a factor \(O(a_j^{-1/2})\) after
integration on its Gaussian scale, in addition to the prefactor
\(a_j^{-1/2}\).  For a left offender, \(e^{-ca_j}\) absorbs the missing local
factor and
\[
        \int_{-\infty}^{-L}\varphi(g)\dd g=e^{-e^L}.
\]
For a right offender, \(e^{-ce^{L/(2\gamma)}a_j}\) absorbs the missing local
factor and
\[
        \int_L^\infty\varphi(g)\dd g
        =1-e^{-e^{-L}}\le e^{-L}.
\]
At least one offender is present.  Summing over the finitely many offender sets
proves \eqref{eq:endpoint-noncompact-localization}.
\end{proof}

\begin{lemma}[Endpoint covariance stability on the saddle scale]
\label{lem:endpoint-covariance-stability}
Let
\[
        H_a=\diag(a_1^{-1/2},\ldots,a_r^{-1/2}),
        \qquad
        \Sigma_*=\Sigma(g^*).
\]
For every fixed \(R<\infty\),
\[
\sup_{\|u\|\le R}
\left\|
D_a^{-1/2}
\bigl[\Sigma(g^*+H_au)-\Sigma_*\bigr]
D_a^{-1/2}
\right\|_{\mathrm{op}}
\longrightarrow0
\tag{8.29}\label{eq:endpoint-covariance-stability}
\]
uniformly over \(\mathfrak E_N(r,\eta)\).  Consequently, uniformly on the
same sets,
\[
        \frac{\det\Sigma(g^*+H_au)}{\det\Sigma_*}\longrightarrow1
\tag{8.30}\label{eq:endpoint-determinant-stability}
\]
and
\[
\left\|
D_a^{1/2}
\bigl[\Sigma(g^*+H_au)^{-1}-\Sigma_*^{-1}\bigr]
D_a^{1/2}
\right\|_{\mathrm{op}}
\longrightarrow0.
\tag{8.31}\label{eq:endpoint-inverse-stability}
\]
\end{lemma}

\begin{proof}
For \(j\le k\), put
\[
        p_{b,j}(g_j)=
        \exp\left(-e^{-g_j}
        \frac{\lambda_{n,b}}{\lambda_{n,a_j}}\right).
\]
Since the threshold events are nested in the chamber,
\[
        \Sigma_{jk}(g)
        =\sum_{b\notin B}p_{b,j}(g_j)(1-p_{b,k}(g_k)).
\tag{8.32}\label{eq:endpoint-covariance-entry}
\]
Lemma \ref{lem:endpoint-rv-sums}, with \(m=1\), gives uniformly on a fixed
saddle neighborhood
\[
        |\partial_{g_j}\Sigma_{jk}(g)|\le Ca_j,
        \qquad
        |\partial_{g_k}\Sigma_{jk}(g)|\le Ca_k.
\tag{8.33}\label{eq:endpoint-covariance-derivatives}
\]
The same estimate, with the evident interpretation, holds on the diagonal.
Hence, for \(\|u\|\le R\),
\[
        |\Sigma_{jk}(g^*+H_au)-\Sigma_{jk}(g^*)|
        \le C_R(\sqrt{a_j}+\sqrt{a_k}).
\]
After division by \(\sqrt{a_ja_k}\), this is at most
\(C_R(a_j^{-1/2}+a_k^{-1/2})\le2C_Ra_1^{-1/2}\).  Since \(r\) is fixed,
\eqref{eq:endpoint-covariance-stability} follows.  By
\eqref{eq:endpoint-covariance-comparison}, the normalized matrices
\(D_a^{-1/2}\Sigma(g)D_a^{-1/2}\) remain in a compact subset of the positive
definite matrices.  Continuity of determinant and inversion on this set gives
\eqref{eq:endpoint-determinant-stability} and
\eqref{eq:endpoint-inverse-stability}.
\end{proof}

\begin{proposition}[Separated endpoint local estimate]
\label{prop:endpoint-rv-joint}
Assume Definition \ref{def:endpoint-rv}.  Fix \(r\ge1\) and \(\eta>0\).
Then
\[
\lim_{N\to\infty}
\sup_{(n,a)\in\mathfrak E_N(r,\eta)}
\left|
\frac{
\Pp(I_{n,a_1}=\cdots=I_{n,a_r}=1)
}{c_\gamma^r/(a_1\cdots a_r)}-1
\right|=0.
\tag{8.34}\label{eq:endpoint-separated-joint}
\]
\end{proposition}

\begin{proof}
On the joint fixed-point event, the specified priorities are necessarily
ordered as
\(y_1<\cdots<y_r\).  Conditional on their centered Gumbels \(g\) in this
chamber, exactly \(j-1\) specified priorities lie below \(y_j\), so the event
is equivalent to \(Z_j(g)=a_j-j\), \(1\le j\le r\).  Therefore
\[
\begin{aligned}
&\Pp(I_{n,a_1}=\cdots=I_{n,a_r}=1)\\
&\quad=\int_{\mathscr C_n}
\prod_{j=1}^r\varphi(g_j)\,
\Pp(Z(g)=z)\dd g.
\end{aligned}
\tag{8.35}\label{eq:endpoint-exact-integral}
\]
Lemma \ref{lem:endpoint-gumbel-saddle} puts a fixed product neighborhood of
\(g^*\) inside the chamber.  Lemma
\ref{lem:endpoint-gumbel-localization} removes the noncompact region on the
scale \(\prod_j a_j^{-1}\), and its compact bound makes the part outside any
fixed saddle neighborhood exponentially negligible on that scale.

Set
\[
        D_n=\diag(M_1'(g_1^*),\ldots,M_r'(g_r^*)),
        \qquad g=g^*+H_au.
\]
For bounded \(u\), Taylor's theorem and
\eqref{eq:endpoint-mean-derivative-bounds} give
\[
        M(g)-z=D_nH_au+r_n(u),
        \qquad
        \|D_a^{-1/2}r_n(u)\|\longrightarrow0
\tag{8.36}\label{eq:endpoint-mean-stability}
\]
uniformly on fixed \(u\)-balls.  Indeed, the \(j\)-th remainder is
\(O_R(1)\), hence is \(o(\sqrt{a_j})\).  The specified density satisfies
\[
        \prod_{j=1}^r\varphi(g_j^*+u_j/\sqrt{a_j})
        =\prod_{j=1}^r\varphi(g_j^*)(1+o(1))
\tag{8.37}\label{eq:endpoint-density-stability}
\]
uniformly on fixed balls.

Apply Lemma \ref{lem:endpoint-gumbel-llt}, then use Lemma
\ref{lem:endpoint-covariance-stability} to replace \(\Sigma(g)\) by
\(\Sigma_*\).  On every fixed \(u\)-ball, the integrand in
\eqref{eq:endpoint-exact-integral}, including the Jacobian
\(|\det H_a|\), is
\[
\begin{aligned}
&(1+o(1))\prod_{j=1}^r\varphi(g_j^*)
(2\pi)^{-r/2}(\det\Sigma_*)^{-1/2}|\det H_a|\\
&\quad\times
\exp\left[-\frac12u^TH_aD_n\Sigma_*^{-1}D_nH_au\right]\dd u.
\end{aligned}
\tag{8.38}\label{eq:endpoint-saddle-integrand}
\]
The matrices in the exponent are uniformly positive definite and uniformly
bounded.  The compact localization estimate
\eqref{eq:endpoint-compact-localization}, after the change of variables,
provides the integrable majorant
\[
        C\left(\prod_{j=1}^ra_j^{-1}\right)
        \exp\left[-c\sum_{j=1}^r\min\{u_j^2,a_j\}\right].
\]
Thus the fixed-ball limit may be extended uniformly to \(\mathbb R^r\).
The Gaussian determinant cancels exactly:
\[
\begin{aligned}
&(2\pi)^{-r/2}(\det\Sigma_*)^{-1/2}|\det H_a|
\int_{\mathbb R^r}
 e^{-\frac12u^TH_aD_n\Sigma_*^{-1}D_nH_au}\dd u\\
&\qquad=|\det D_n|^{-1}.
\end{aligned}
\tag{8.39}\label{eq:endpoint-determinant-cancellation}
\]
Consequently,
\[
\Pp(I_{n,a_1}=\cdots=I_{n,a_r}=1)
=(1+o(1))\prod_{j=1}^r
\frac{\varphi(g_j^*)}{M_j'(g_j^*)}.
\tag{8.40}\label{eq:endpoint-saddle-product}
\]
By \eqref{eq:endpoint-saddle-asymptotics},
\[
        \varphi(g_j^*)\longrightarrow
        K_\gamma e^{-K_\gamma},
        \qquad
        M_j'(g_j^*)\sim a_j/\gamma.
\]
Each factor in \eqref{eq:endpoint-saddle-product} is therefore asymptotic to
\(c_\gamma/a_j\), uniformly over the admissible family.  This proves
\eqref{eq:endpoint-separated-joint}.
\end{proof}


\begin{theorem}[Endpoint fixed-point universality]
\label{thm:endpoint-fixed-universality}
Assume the terminal rate array is uniformly endpoint regularly varying with
index \(\gamma>0\) in the sense of Definition \ref{def:endpoint-rv}.  Let
\[
        \Xi_n=\sum_{a=1}^n I_{n,a}\delta_{\log(a/b_n)},
\]
where \(b_n\to\infty\) and \(b_n=o(n)\).  For every bounded interval
\(J\subset\mathbb R\), the restriction \(\Xi_n|_J\), regarded as a random
locally finite counting measure with the vague topology, satisfies
\[
        \Xi_n|_J\Rightarrow \PPP(c_\gamma\dd x)|_J.
\tag{8.41}\label{eq:endpoint-ppp-limit}
\]
Equivalently, counts in finitely many disjoint bounded intervals converge
jointly to independent Poisson random variables with means \(c_\gamma\) times
their lengths.  Moreover, for every fixed \(r\ge1\),
\[
        \Ee(F_n)_r\sim(c_\gamma\log n)^r.
\tag{8.42}\label{eq:endpoint-factorial-moments}
\]
In particular,
\[
        \Ee C_1(\sigma_n)\sim c_\gamma\log n,
        \qquad
        \frac{C_1(\sigma_n)}{c_\gamma\log n}
        \longrightarrow1
        \quad\text{in probability}.
\tag{8.43}\label{eq:endpoint-fixed-lln}
\]
\end{theorem}

\begin{proof}
We first prove the local point-process assertion.  For every bounded interval
\(I=(u,v]\), the elementary logarithmic Riemann sum gives
\[
        \sum_{\log(a/b_n)\in I}\frac1a\longrightarrow v-u=|I|.
\tag{8.44}\label{eq:endpoint-harmonic-window}
\]
All coordinates in this sum tend to infinity and are \(o(n)\).  Fix disjoint
bounded intervals \(I_1,\ldots,I_m\) and nonnegative integers
\(k_1,\ldots,k_m\), and put \(k=\sum_\ell k_\ell\).  Expand
\[
        \Ee\prod_{\ell=1}^m(\Xi_n(I_\ell))_{k_\ell}
\]
over ordered distinct terminal coordinates, colored according to their
intervals.  Restrict first to tuples whose logarithms are pairwise more than
\(\eta\) apart.  After sorting the coordinates, Proposition
\ref{prop:endpoint-rv-joint} applies with ratio-separation parameter
\(e^\eta-1\), uniformly over the entire restricted sum.  Its contribution is
therefore
\[
        c_\gamma^k(1+o(1))
        \sum_{\eta\text{-separated}}^*
        \frac1{a_1\cdots a_k}.
\tag{8.45}\label{eq:endpoint-separated-window-sum}
\]

The logarithmic diagonal neighborhood has negligible harmonic weight as
\(\eta\downarrow0\).  Indeed, uniformly for \(a\) in a fixed multiplicative
window,
\[
        \sum_{b:\,|\log b-\log a|\le\eta}\frac1b
        \le C\eta+o(1),
\tag{8.46}\label{eq:endpoint-log-diagonal-bound}
\]
and hence a union bound over pairs gives an \(O_{k,I_1,\ldots,I_m}(\eta)+o(1)\)
bound for the weighted sum of excluded tuples.  Lemma
\ref{lem:endpoint-rv-product-bound} gives the same bound, up to a fixed
constant, for their probability contribution.  Using
\eqref{eq:endpoint-harmonic-window} in
\eqref{eq:endpoint-separated-window-sum} and then letting
\(\eta\downarrow0\), we obtain
\[
        \Ee\prod_{\ell=1}^m(\Xi_n(I_\ell))_{k_\ell}
        \longrightarrow
        \prod_{\ell=1}^m(c_\gamma|I_\ell|)^{k_\ell}.
\tag{8.47}\label{eq:endpoint-mixed-window-factorials}
\]
These are the mixed factorial moments of independent Poisson variables, so the
interval-count vectors converge to the asserted independent Poisson vector.

We now upgrade count convergence to convergence of point measures.  By the
one-point case of Lemma \ref{lem:endpoint-rv-product-bound},
\[
        \sup_n\Ee\Xi_n(J)
        \le C\sup_n\sum_{\log(a/b_n)\in J}\frac1a<\infty
\tag{8.48}\label{eq:endpoint-ppp-tightness}
\]
for every bounded \(J\), after ignoring finitely many initial \(n\).  Let
\(f\ge0\) be continuous with compact support in \(J\).  Choose a finite
partition of \(J\) into half-open intervals and a step function \(f_\pi\),
constant on the partition cells, with
\(\|f-f_\pi\|_\infty\to0\) as the mesh tends to zero.  Since
\(|e^{-x}-e^{-y}|\le|x-y|\) for \(x,y\ge0\),
\[
\left|
\Ee e^{-\langle\Xi_n,f\rangle}
-\Ee e^{-\langle\Xi_n,f_\pi\rangle}
\right|
\le\|f-f_\pi\|_\infty\Ee\Xi_n(J),
\tag{8.49}\label{eq:endpoint-laplace-step-error}
\]
uniformly in \(n\).  For fixed \(\pi\), the interval-count convergence just
proved gives
\[
\Ee e^{-\langle\Xi_n,f_\pi\rangle}
\longrightarrow
\exp\left[-c_\gamma\int_J(1-e^{-f_\pi(x)})\dd x\right].
\]
Letting the mesh tend to zero in view of
\eqref{eq:endpoint-laplace-step-error} yields
\[
\Ee e^{-\langle\Xi_n,f\rangle}
\longrightarrow
\exp\left[-c_\gamma\int_J(1-e^{-f(x)})\dd x\right],
\tag{8.50}\label{eq:endpoint-laplace-functional}
\]
the Laplace functional of \(\PPP(c_\gamma\dd x)|_J\).  Laplace functionals
on nonnegative compactly supported continuous functions determine the law of a
locally finite counting measure and are convergence determining for the vague
topology.  This proves \eqref{eq:endpoint-ppp-limit}; the uniform first-moment
bound \eqref{eq:endpoint-ppp-tightness} also gives tightness directly.

It remains to prove the full factorial moments.  Put
\[
        L_n=\lceil\log n\rceil,
        \qquad
        U_n=\lfloor n/\log n\rfloor,
        \qquad
        H_n=\sum_{L_n\le a\le U_n}\frac1a.
\]
Then \(H_n\sim\log n\).  Fix once and for all a positive separation
\(\eta_0\).  The harmonic weight of ordered \(r\)-tuples in
\([L_n,U_n]\) containing a pair with
\(|\log a_i-\log a_j|\le\eta_0\) is \(O_r(H_n^{r-1})\).  To see this, fix
one coordinate \(a_i\), use
\[
        \sum_{b:\,|\log b-\log a_i|\le\eta_0}\frac1b
        \le C_{\eta_0},
\]
and sum the remaining \(r-2\) coordinates freely.  Repeated coordinates are
included in this bound.  Therefore the separated weighted sum is
\[
        H_n^r+O_r(H_n^{r-1}).
\tag{8.51}\label{eq:endpoint-global-separated-weight}
\]
On the separated tuples, Proposition \ref{prop:endpoint-rv-joint} applies
uniformly and gives the contribution
\[
        c_\gamma^rH_n^r+o(H_n^r).
\]
The nonseparated probability contribution is \(O_r(H_n^{r-1})\) by Lemma
\ref{lem:endpoint-rv-product-bound} and the preceding harmonic estimate.
Thus, for
\(G_n=\sum_{L_n\le a\le U_n}I_{n,a}\),
\[
        \Ee(G_n)_r\sim(c_\gamma\log n)^r.
\tag{8.52}\label{eq:endpoint-truncated-factorials}
\]

Finally, the harmonic mass outside \([L_n,U_n]\) is
\[
        H_n^{\mathrm{out}}
        :=\sum_{a<L_n}\frac1a+
          \sum_{a>U_n}\frac1a
        =O(\log\log n).
\]
By the product envelope, the total contribution to \(\Ee(F_n)_r\) from
ordered tuples containing at least one omitted coordinate is at most
\[
        C_rH_n^{\mathrm{out}}
        \left(\sum_{a=1}^n\frac1a\right)^{r-1}
        =O_r((\log n)^{r-1}\log\log n)
        =o((\log n)^r).
\]
Combining this with \eqref{eq:endpoint-truncated-factorials} proves
\eqref{eq:endpoint-factorial-moments}.

Taking \(r=1,2\) gives
\(\Ee F_n\sim c_\gamma\log n\) and
\(\Var(F_n)=o((\log n)^2)\).  Chebyshev's inequality yields
\eqref{eq:endpoint-fixed-lln}.
\end{proof}


\subsection{Further endpoint limits}

The same regular-variation calculation also gives two boundary limits that are
useful for interpreting the endpoint process.  The first is a mesoscopic
one-coordinate displacement law; the second is a microscopic infinite random
permutation.

\begin{theorem}[Endpoint displacement law]\label{thm:endpoint-displacement}
Assume Definition \ref{def:endpoint-rv}.  Let
$a_{n,1},\ldots,a_{n,m}$ be distinct terminal coordinates, where $m$ is fixed,
$a_{n,j}\to\infty$, and $a_{n,j}=o(n)$ for every $j$.  Then
\[
        \left(\frac{\tau_n(a_{n,1})}{a_{n,1}},\ldots,
        \frac{\tau_n(a_{n,m})}{a_{n,m}}\right)
        \Rightarrow
        \left(A_\gamma E_1^{-1/\gamma},\ldots,
        A_\gamma E_m^{-1/\gamma}\right),
\]
where $E_1,\ldots,E_m$ are independent $\Exp(1)$ variables and
$A_\gamma=\Gamma(1+1/\gamma)$.  Equivalently,
\[
        \log\frac{\tau_n(a_{n,j})}{a_{n,j}}
        \Rightarrow X_\gamma,
\]
where $X_\gamma$ has density
\[
        k_\gamma(x)=\gamma K_\gamma e^{-\gamma x}
        \exp\{-K_\gamma e^{-\gamma x}\},
        \qquad K_\gamma=A_\gamma^\gamma.
\]
\end{theorem}

\begin{proof}
Write
\[
        Y_{n,a}=\log\lambda_{n,a}+G_a,
\]
with independent standard Gumbels.  Conditional on $G_a=g$,
\[
        \tau_n(a)-1
        =\sum_{b\ne a}\one_{\{Y_{n,b}<Y_{n,a}\}}
\]
is a sum of independent Bernoulli variables with mean
\[
        M_{n,a}(g)=\sum_{b\ne a}
        \exp\left(-e^{-g}\frac{\lambda_{n,b}}{\lambda_{n,a}}\right).
\]
Lemma \ref{lem:endpoint-rv-sums} gives, uniformly for $g$ in compact sets,
\[
        \frac{M_{n,a}(g)}a\longrightarrow
        A_\gamma e^{g/\gamma}.
\]
The conditional variance is at most the conditional mean, hence is $O(a)$ on
compact $g$-sets.  Chebyshev's inequality therefore gives, uniformly on such
sets,
\[
        \frac{\tau_n(a)}a-A_\gamma e^{G_a/\gamma}
        \longrightarrow0
        \quad\text{in probability}.
\]
Truncating the Gumbel variables to a fixed compact set and then letting the
truncation grow proves the one-coordinate assertion.  For finitely many
specified coordinates, delete them from the Bernoulli sums.  Lemma
\ref{lem:endpoint-rv-sums} is unchanged after finitely many deletions, and a
union bound over the conditional Chebyshev estimates proves the joint
convergence.  The limiting variables are independent because the specified
Gumbels are independent.

Finally, if $E=e^{-G}\sim\Exp(1)$, then
$A_\gamma e^{G/\gamma}=A_\gamma E^{-1/\gamma}$.  The change of variables
$x=\log A_\gamma+G/\gamma$ gives the displayed density $k_\gamma$.
\end{proof}

\subsection{Mesoscopic endpoint tangent measure}
\label{subsec:endpoint-tangent}

The microscopic limit in Theorem \ref{thm:microscopic-endpoint-permutation} is
an infinite random permutation.  On every mesoscopic scale containing
$b_n\to\infty$ terminal coordinates, the correct object is instead a dense
normalized graph measure.  The unnormalized graph cannot have a locally finite
point-process limit, because a compact source window contains order $b_n$
points.

Define
\[
        \mathcal K_\gamma(x,y)
        =\frac1y k_\gamma\!\left(\log\frac yx\right)
        =\gamma K_\gamma\frac{x^\gamma}{y^{\gamma+1}}
        \exp\left\{-K_\gamma\left(\frac xy\right)^\gamma\right\},
        \qquad x,y>0.
\]

\begin{theorem}[Mesoscopic endpoint tangent measure]
\label{thm:endpoint-tangent-measure}
Assume Definition \ref{def:endpoint-rv}.  Let $b_n\to\infty$ with
$b_n=o(n)$, and put
\[
        \mathcal M_n^{(b_n)}
        =\frac1{b_n}\sum_{a=1}^n
        \delta_{(a/b_n,\tau_n(a)/b_n)}.
\]
Each \(\mathcal M_n^{(b_n)}\) is a Radon measure on compact subsets of
\((0,\infty)^2\), and the density \(\mathcal K_\gamma\) is locally
integrable there.  Then, vaguely in probability on $(0,\infty)^2$,
\[
        \mathcal M_n^{(b_n)}
        \longrightarrow
        \mathcal K_\gamma(x,y)\dd x\dd y.
\]
More precisely, for every $\varphi\in C_c((0,\infty)^2)$,
\[
\begin{aligned}
        \Ee\left|
        \langle\mathcal M_n^{(b_n)},\varphi\rangle
        -\int_0^\infty\int_0^\infty
        \varphi(x,y)\mathcal K_\gamma(x,y)\dd y\dd x
        \right|\longrightarrow0.
\end{aligned}
\]
The limiting kernel is bistochastic with respect to Lebesgue measure.
\end{theorem}

\begin{lemma}[Uniform mesoscopic Laplace sums]
\label{lem:mesoscopic-laplace-sums}
Fix $0<x_-<x_+<\infty$ and $0<u_-<u_+<\infty$.  Uniformly for
$x_-b_n\le a\le x_+b_n$ and $u\in[u_-,u_+]$,
\[
        \frac1{b_n}\sum_{c=1}^n
        \exp\left\{-u\frac{\lambda_{n,c}}{\lambda_{n,a}}\right\}
        =\frac a{b_n}A_\gamma u^{-1/\gamma}+o(1).
\]
For every fixed $u_->0$, the left-hand side is uniformly bounded for all
$u\ge u_-$.
\end{lemma}

\begin{proof}
Put $X_n=n/b_n\to\infty$ and choose
\[
        R_n=(\log(e+X_n))^{4/\gamma},
        \qquad B_n=\lceil2x_+b_nR_n\rceil.
\]
Then $B_n=o(n)$.  Uniform endpoint regular variation applies to all
$a\in[x_-b_n,x_+b_n]$ and $c\le B_n$.  On every fixed multiplicative interval
$\delta a\le c\le Ra$, the uniform convergence theorem for slowly varying
functions gives
\[
        \lambda_{n,c}/\lambda_{n,a}=(c/a)^\gamma(1+o(1))
\]
uniformly.  Hence
\[
\begin{aligned}
        \frac1{b_n}\sum_{\delta a\le c\le Ra}
        e^{-u\lambda_{n,c}/\lambda_{n,a}}
        =\frac a{b_n}\int_\delta^R e^{-uv^\gamma}\dd v+o(1).
\end{aligned}
\]
The range $c<\delta a$ contributes at most $x_+\delta$.  Potter's bound with
exponent $\gamma/2$ gives, for $Ra<c\le B_n$,
\[
        \lambda_{n,c}/\lambda_{n,a}
        \ge c(c/a)^{\gamma/2},
\]
so the normalized contribution is bounded by
\[
        C\frac a{b_n}\int_R^\infty e^{-cu_-v^{\gamma/2}}\dd v.
\]
For $c>B_n$, monotonicity and Potter's bound give
\[
        \lambda_{n,c}/\lambda_{n,a}\ge cR_n^{\gamma/2},
\]
and therefore
\[
        \frac1{b_n}\sum_{c>B_n}
        e^{-u\lambda_{n,c}/\lambda_{n,a}}
        \le X_n\exp\{-cu_-R_n^{\gamma/2}\}=o(1).
\]
Let first $n\to\infty$, then $\delta\downarrow0$ and $R\uparrow\infty$.
Since
\[
        \int_0^\infty e^{-uv^\gamma}\dd v
        =A_\gamma u^{-1/\gamma},
\]
the first assertion follows.  Monotonicity in $u$ gives the final bound.
\end{proof}

Let $E_{n,a}=\lambda_{n,a}T_{n,a}$.  These variables are independent and
standard exponential.  Conditional on $E_{n,a}=u$,
\[
        \tau_n(a)=1+\sum_{c\ne a}\xi_{n,a,c}(u),
        \qquad
        \Pp(\xi_{n,a,c}(u)=1)
        =e^{-u\lambda_{n,c}/\lambda_{n,a}},
\]
with independent summands.  Put
\[
        r_{n,a}(u)=1+\sum_{c\ne a}
        e^{-u\lambda_{n,c}/\lambda_{n,a}}.
\]

\begin{lemma}[Conditional rank self-averaging]
\label{lem:conditional-rank-self-averaging}
Fix $0<x_-<x_+<\infty$.  For every $\varepsilon,\delta>0$,
\[
\begin{aligned}
        \sup_{x_-b_n\le a\le x_+b_n}
        \Pp\bigl(|\tau_n(a)-r_{n,a}(E_{n,a})|>\delta b_n\bigr)
        \le\varepsilon+\frac{C_\varepsilon}{\delta^2b_n}
\end{aligned}
\]
for all large $n$.  Moreover, uniformly for
$a/b_n\in[x_-,x_+]$ and $u$ in a compact subset of $(0,\infty)$,
\[
        \frac{r_{n,a}(u)}{b_n}
        \longrightarrow
        \frac a{b_n}A_\gamma u^{-1/\gamma}.
\]
\end{lemma}

\begin{proof}
Conditionally on $E_{n,a}=u$,
\[
        \Var(\tau_n(a)\mid E_{n,a}=u)\le r_{n,a}(u).
\]
On $\{E_{n,a}\ge\varepsilon\}$, Lemma
\ref{lem:mesoscopic-laplace-sums} and monotonicity in $u$ give
$r_{n,a}(E_{n,a})\le C_\varepsilon b_n$.  Conditional Chebyshev therefore
bounds the deviation probability on this event by
$C_\varepsilon/(\delta^2b_n)$.  The complementary event has probability at
most $\varepsilon$.  The second assertion follows immediately from Lemma
\ref{lem:mesoscopic-laplace-sums}; omitting the $c=a$ summand and adding the
leading one changes the sum by $O(1)$.
\end{proof}

\begin{proof}[Proof of Theorem \ref{thm:endpoint-tangent-measure}]
Fix $\varphi\in C_c((0,\infty)^2)$, and choose
$0<x_-<x_+<\infty$ containing the first-coordinate projection of its support.
Set
\[
        Z_n(\varphi)=\frac1{b_n}\sum_a
        \varphi(a/b_n,\tau_n(a)/b_n)
\]
and
\[
        \widehat Z_n(\varphi)=\frac1{b_n}\sum_a
        \varphi(a/b_n,r_{n,a}(E_{n,a})/b_n).
\]
Only $O(b_n)$ summands are nonzero.  Uniform continuity of $\varphi$ and Lemma
\ref{lem:conditional-rank-self-averaging} give
\[
        \Ee|Z_n(\varphi)-\widehat Z_n(\varphi)|\longrightarrow0.
\]
Indeed, split according as the two second coordinates differ by at most
$\delta$, use the modulus of continuity on that event, and then let
$n\to\infty$, $\delta\downarrow0$, and $\varepsilon\downarrow0$.

The summands in $\widehat Z_n(\varphi)$ are independent functions of the
independent variables $E_{n,a}$.  Hence
\[
        \Var(\widehat Z_n(\varphi))\le C_\varphi/b_n\longrightarrow0.
\]
It remains to compute the expectations.  If $U\sim\Exp(1)$, Lemma
\ref{lem:conditional-rank-self-averaging}, uniformly on compact $u$-sets,
gives
\[
\begin{aligned}
        \Ee\widehat Z_n(\varphi)
        &\longrightarrow
        \int_0^\infty\int_0^\infty
        \varphi(x,xA_\gamma u^{-1/\gamma})e^{-u}\dd u\dd x.
\end{aligned}
\]
To justify the passage, truncate $u$ to $[\varepsilon,M]$, use uniform
convergence and an ordinary Riemann sum, and then use the exponential mass of
$(0,\varepsilon)\cup(M,\infty)$.

For fixed $x$, make the change of variables
$y=xA_\gamma u^{-1/\gamma}$.  Then
$u=K_\gamma(x/y)^\gamma$ and
\[
        e^{-u}\left|\frac{\dd u}{\dd y}\right|
        =\mathcal K_\gamma(x,y).
\]
This proves the $L^1$ convergence of test-function integrals.  A countable
convergence-determining subset of $C_c((0,\infty)^2)$ then gives vague
convergence in probability.

Finally, the first marginal identity follows from the same change of variables:
$\mathcal K_\gamma(x,\cdot)$ is the density of
$xA_\gamma U^{-1/\gamma}$.  For the second marginal, put
$v=K_\gamma(x/y)^\gamma$ to obtain
\[
\begin{aligned}
        \int_0^\infty\mathcal K_\gamma(x,y)\dd x
        &=K_\gamma^{-1/\gamma}
        \int_0^\infty v^{1/\gamma}e^{-v}\dd v=1.
\end{aligned}
\]
Thus the kernel is bistochastic.
\end{proof}

\begin{corollary}[Rectangle law]
\label{cor:endpoint-tangent-rectangles}
For bounded intervals $I,J\Subset(0,\infty)$,
\[
        \frac1{b_n}\#\left\{a:
        a/b_n\in I,\ \tau_n(a)/b_n\in J\right\}
        \xrightarrow{\Pp}
        \int_I\int_J\mathcal K_\gamma(x,y)\dd y\dd x.
\]
\end{corollary}

\begin{proof}
The boundary of $I\times J$ has zero mass under the continuous limiting
density, so the assertion follows from vague convergence and approximation by
compactly supported continuous functions.
\end{proof}

\begin{proposition}[Multiplicative semigroup and diagonal traces]
\label{prop:endpoint-kernel-powers}
For every $r\ge1$,
\[
        \mathcal K_\gamma^{\,r}(x,y)
        =\frac1y k_\gamma^{*r}\!\left(\log\frac yx\right).
\]
In particular,
\[
        \mathcal K_\gamma^{\,r}(x,x)
        =\frac1x k_\gamma^{*r}(0).
\]
\end{proposition}

\begin{proof}
For $r=2$, substitute $z=xe^w$ in the composition integral:
\[
\begin{aligned}
        \int_0^\infty\mathcal K_\gamma(x,z)
        \mathcal K_\gamma(z,y)\dd z
        &=\frac1y\int_{\mathbb R}k_\gamma(w)
        k_\gamma\!\left(\log\frac yx-w\right)\dd w.
\end{aligned}
\]
This is $y^{-1}k_\gamma^{*2}(\log(y/x))$.  Iterate.
\end{proof}


\subsection{Microscopic endpoint permutation}\label{subsec:microscopic-endpoint}

\begin{theorem}[Microscopic endpoint permutation]
\label{thm:microscopic-endpoint-permutation}
Assume Definition \ref{def:endpoint-rv}.  Extend the finite permutation
\(\tau_n\in S_n\) to a permutation \(\bar\tau_n\) of \(\mathbb N\) by
\[
        \bar\tau_n(a)=
        \begin{cases}
        \tau_n(a),&a\le n,\\
        a,&a>n.
        \end{cases}
\]
Then, in the product topology on \(\mathbb N^{\mathbb N}\),
\[
        \bar\tau_n\Rightarrow\tau^{(\gamma)},
\tag{8.53}\label{eq:microscopic-endpoint-convergence}
\]
where, for independent \(E_1,E_2,\ldots\sim\Exp(1)\),
\[
        \tau^{(\gamma)}(a)
        =1+\#\left\{b\ne a:
        \frac{E_b}{b^\gamma}>\frac{E_a}{a^\gamma}\right\}.
\]
Almost surely, \(\tau^{(\gamma)}\) is a permutation of \(\mathbb N\).
\end{theorem}

\begin{proof}
Normalize by the first terminal rate and put
\[
        r_{n,b}=\frac{\lambda_{n,b}}{\lambda_{n,1}},
        \qquad
        Q_{n,b}=\lambda_{n,1}T_{n,b}=\frac{E_b}{r_{n,b}},
\]
using independent \(E_b\sim\Exp(1)\).  For every fixed \(B\), uniform
endpoint regular variation and the uniform convergence theorem for slowly
varying functions give
\[
        (r_{n,1},\ldots,r_{n,B})
        \longrightarrow(1,2^\gamma,\ldots,B^\gamma).
\]
Hence, under this coupling,
\[
        (Q_{n,1},\ldots,Q_{n,B})
        \longrightarrow
        (E_1,E_2/2^\gamma,\ldots,E_B/B^\gamma)
\]
almost surely.

We remove the tail uniformly.  Fix \(m\) and \(\varepsilon>0\).  For fixed
\(a\le m\), the ratios \(r_{n,a}\) remain bounded, so
\[
        \Pp(Q_{n,a}<\varepsilon)\le C_m\varepsilon.
\]
Choose \(W_n=\lfloor n/\log n\rfloor\).  Potter's bound, applied with
reference coordinate \(1\), gives
\[
        r_{n,b}\ge cb^{\gamma/2},
        \qquad B<b\le W_n,
\]
for all sufficiently large \(n\).  By monotonicity, the same lower bound at
\(b=W_n\) applies to every \(b>W_n\).  Consequently,
\[
\begin{aligned}
        \sum_{b>B}\Pp(Q_{n,b}>\varepsilon)
        &=\sum_{b>B}e^{-\varepsilon r_{n,b}}\\
        &\le\sum_{b>B}e^{-c\varepsilon b^{\gamma/2}}
        +n e^{-c\varepsilon W_n^{\gamma/2}}.
\end{aligned}
\tag{8.54}\label{eq:microscopic-tail-truncation}
\]
The second term tends to zero, and the first tends to zero as \(B\to\infty\),
uniformly for all sufficiently large \(n\).  Thus, outside an event of
probability at most \(C_m\varepsilon+o_B(1)\), every comparison determining
the ranks of the first \(m\) coordinates is decided by the first \(B\) clocks.
The same argument applies to the limiting clocks \(E_b/b^\gamma\).  Finite
collections of all these clocks have no ties.  Taking first \(n\to\infty\),
then \(B\to\infty\), and finally \(\varepsilon\downarrow0\), proves
finite-dimensional convergence of
\((\bar\tau_n(1),\ldots,\bar\tau_n(m))\) for every fixed \(m\), which is
\eqref{eq:microscopic-endpoint-convergence}.

For every \(t>0\),
\[
        \sum_{b=1}^\infty
        \Pp(E_b/b^\gamma>t)
        =\sum_{b=1}^\infty e^{-b^\gamma t}<\infty.
\]
Borel--Cantelli implies that only finitely many limiting clocks exceed any
positive threshold.  Moreover,
\[
        \Pp\left(\sup_{b\ge1}\frac{E_b}{b^\gamma}>t\right)
        \le\sum_{b\ge1}e^{-b^\gamma t}\longrightarrow0
        \qquad(t\to\infty),
\]
so the clock set has a finite maximum.  The clocks are almost surely distinct,
and zero is their only accumulation point.  Their decreasing enumeration
therefore contains every label exactly once and assigns every positive integer
rank exactly once.  Thus \(\tau^{(\gamma)}\) is almost surely a permutation of
\(\mathbb N\).
\end{proof}


\begin{remark}[The Gumbel origin of the constant]\label{rem:gumbel-endpoint-constant}
The preceding proof makes the constant transparent.  Conditional on
\(Y_{n,a}=\log\lambda_{n,a}+g\), the expected number of lower terminal priorities
is
\[
        \sum_b\exp\left(-e^{-g}\frac{\lambda_{n,b}}{\lambda_{n,a}}\right)
        \sim aA_\gamma e^{g/\gamma}.
\]
The fixed-point saddle is therefore
\(A_\gamma e^{g/\gamma}=1\), so
\(e^{-g}=K_\gamma=A_\gamma^\gamma\).  The density of the specified centered
Gumbel at the saddle is \(K_\gamma e^{-K_\gamma}\), while the derivative of the
mean with respect to \(g\) is \(a/\gamma\).  Their quotient is exactly
\(c_\gamma/a\), with
\[
        c_\gamma=\gamma K_\gamma e^{-K_\gamma}.
\]
For \(\gamma=1\), \(A_\gamma=K_\gamma=1\), and the constant reduces to the
Sukhatme value \(e^{-1}\).
\end{remark}

% END ARCHIVED BLOCK: endpoint-regular-universality
% BEGIN ARCHIVED BLOCK: endpoint-short-cycles
\section{Endpoint short cycles}
\label{sec:endpoint-short-cycles}

Recall that the terminal-coordinate permutation $\tau_n$ is conjugate to
$\sigma_n$ by reversal.  In particular,
\[
        C_r(\tau_n)=C_r(\sigma_n),\qquad r\ge1.
\]
The multiplicative kernel in Proposition
\ref{prop:endpoint-kernel-powers} determines the asymptotic number of every
fixed endpoint cycle length.  Unlike the regular positive-profile regime, each
of these expected cycle counts diverges logarithmically.

For $r\ge1$, put
\[
        \alpha_{\gamma,r}=\frac1r k_\gamma^{*r}(0).
\]
Thus $\alpha_{\gamma,1}=c_\gamma$.

\begin{theorem}[Endpoint short-cycle law]
\label{thm:endpoint-short-cycles}
Assume Definition \ref{def:endpoint-rv}.  For every fixed $R\ge1$ and fixed
nonnegative integers $m_1,\ldots,m_R$,
\[
        \Ee\prod_{r=1}^R(C_r(\sigma_n))_{m_r}
        \sim\prod_{r=1}^R
        (\alpha_{\gamma,r}\log n)^{m_r}.
\]
Consequently, for every fixed $r$,
\[
        \Ee C_r(\sigma_n)
        \sim\frac{k_\gamma^{*r}(0)}r\log n,
\]
and, jointly for every fixed $R$,
\[
        \left(
        \frac{C_1}{\alpha_{\gamma,1}\log n},\ldots,
        \frac{C_R}{\alpha_{\gamma,R}\log n}
        \right)\longrightarrow(1,\ldots,1)
\]
in probability and in every fixed $L^p$, $p<\infty$.
\end{theorem}

The proof has three inputs.  First, we prove a source-to-target endpoint entry
estimate by an endpoint-specific anisotropic local central limit theorem,
mixed noncompact localization under one tilted law, and covariance stability
on the saddle scale.  Second, a death-process envelope localizes every fixed
cycle component in logarithmic coordinates, including components crossing the
mesoscopic cutoff.  Third, a logarithmic trace Riemann sum identifies the
coefficient $k_\gamma^{*r}(0)/r$.

\subsection{General endpoint point constraints}

For fixed $\ell\ge1$, $\eta>0$, and $M<\infty$, define
\[
\begin{aligned}
\mathfrak Q_N(\ell,\eta,M)=\bigl\{&(n,p_1,\ldots,p_\ell,
 q_1,\ldots,q_\ell):
 N\le p_j,q_j\le n/N,\\
& p_1,\ldots,p_\ell\text{ distinct},\quad
 q_1<\cdots<q_\ell,\quad
 q_{j+1}/q_j\ge e^\eta,\\
& e^{-M}\le p_j/q_j\le e^M\quad(1\le j\le\ell)
\bigr\}.
\end{aligned}
\tag{10.1}\label{eq:endpoint-point-admissible}
\]
A uniform assertion below means that the corresponding supremum over
$\mathfrak Q_N(\ell,\eta,M)$ tends to zero as $N\to\infty$.
The coordinates $p_j$ are called \emph{source terminal coordinates}, and the
$q_j$ are the target terminal positions.

Let $P=\{p_1,\ldots,p_\ell\}$ and put
\[
        R_{n,j}=\frac{\lambda_{n,p_j}}{\lambda_{n,q_j}},
        \qquad
        z_j=q_j-j.
\]
For $h=(h_1,\ldots,h_\ell)\in\mathbb R^\ell$, define thresholds
\[
        y_j(h_j)=\log\lambda_{n,q_j}+h_j
\]
and cumulative lower-priority counts
\[
        Z_j(h)=\sum_{b\notin P}
        \one_{\{Y_{n,b}<y_j(h_j)\}},
        \qquad 1\le j\le\ell.
\]
The relevant chamber is
\[
        \mathscr C_{n,q}
        =\{h:y_1(h_1)<\cdots<y_\ell(h_\ell)\}.
\]
Finally, for $R>0$ let
\[
        \varphi_R(h)=R e^{-h}\exp\{-R e^{-h}\}.
\tag{10.2}\label{eq:shifted-gumbel-density}
\]
This is the density of $\log R+G$ when $G$ is standard Gumbel.

\begin{lemma}[Source-to-target endpoint saddle geometry]
\label{lem:endpoint-point-saddle}
Uniformly over $\mathfrak Q_N(\ell,\eta,M)$, the ratios $R_{n,j}$ remain in a
fixed compact subinterval of $(0,\infty)$.  For each $j$, the equation
\[
        M_j(h_j^*)=z_j,
        \qquad
        M_j(h)=\Ee Z_j(h),
\tag{10.3}\label{eq:endpoint-point-saddle-equation}
\]
has a unique solution, and
\[
        h_j^*\longrightarrow g_\gamma:=-\log K_\gamma,
        \qquad
        d_{n,j}:=M_j'(h_j^*)\sim\frac{q_j}{\gamma}.
\tag{10.4}\label{eq:endpoint-point-saddle-asymptotics}
\]
For every fixed $L<\infty$,
\[
        \sup_{|h|\le L}|M_j''(h)|\le C_Lq_j,
        \qquad
        c_Lq_j\le M_j'(h)\le C_Lq_j
\tag{10.5}\label{eq:endpoint-point-mean-derivatives}
\]
whenever $|h-h_j^*|\le1$.  Moreover, there are
$\delta_{\eta,M}>0$, $c_{\eta,M}>0$, and $N_0$ such that, for $N\ge N_0$,
\[
        y_{j+1}(h_{j+1})-y_j(h_j)\ge c_{\eta,M}
\tag{10.6}\label{eq:endpoint-point-chamber-margin}
\]
whenever $|h_j-h_j^*|\le\delta_{\eta,M}$ for every $j$.  Thus this product
neighborhood is contained in $\mathscr C_{n,q}$.
\end{lemma}

\begin{proof}
Uniform regular variation on the fixed multiplicative range
$e^{-M}\le p_j/q_j\le e^M$ gives
\[
        R_{n,j}
        =\left(\frac{p_j}{q_j}\right)^\gamma(1+o(1))
\tag{10.7}\label{eq:endpoint-point-rate-ratio}
\]
uniformly, proving the compactness assertion.  By Lemma
\ref{lem:endpoint-rv-sums}, including its finite-deletion clause, uniformly on
compact $h$-sets,
\[
        M_j(h)=q_jA_\gamma e^{h/\gamma}(1+o(1)),
\tag{10.8}\label{eq:endpoint-point-mean-asymptotic}
\]
\[
        M_j'(h)=\frac{q_jA_\gamma}{\gamma}
        e^{h/\gamma}(1+o(1)).
\tag{10.9}\label{eq:endpoint-point-mean-prime}
\]
The function $M_j$ is continuous and strictly increasing from $0$ to
$n-\ell$, so the saddle exists uniquely.  Since $z_j/q_j\to1$,
\eqref{eq:endpoint-point-mean-asymptotic} gives
$A_\gamma e^{h_j^*/\gamma}\to1$, equivalently
$h_j^*\to-\log K_\gamma$.  Substitution into
\eqref{eq:endpoint-point-mean-prime} gives
$d_{n,j}\sim q_j/\gamma$.

Termwise differentiation expresses $M_j''$ as a linear combination of the
weighted sums in Lemma \ref{lem:endpoint-rv-sums} with powers one and two.
This gives \eqref{eq:endpoint-point-mean-derivatives}.

Finally, the target separation and the Potter lower bound imply that, for some
$c_\eta>0$ and all large $N$,
\[
        \log\frac{\lambda_{n,q_{j+1}}}{\lambda_{n,q_j}}
        \ge2c_\eta.
\tag{10.10}\label{eq:endpoint-point-rate-gap}
\]
Indeed, on a fixed compact range of $q_{j+1}/q_j$ this follows from the uniform
convergence theorem, while larger ratios are handled by
$\lambda_{n,q_{j+1}}/\lambda_{n,q_j}
 \ge c(q_{j+1}/q_j)^{\gamma/2}$.
Since all $h_j^*$ converge to the same constant, a sufficiently small fixed
$\delta_{\eta,M}$ gives \eqref{eq:endpoint-point-chamber-margin}.
\end{proof}

\begin{lemma}[Source-to-target endpoint anisotropic local CLT]
\label{lem:endpoint-point-llt}
Let
\[
        D_q=\diag(q_1,\ldots,q_\ell),
        \qquad
        \Sigma(h)=\Var Z(h),
\]
and let $\mathcal U_n$ be the saddle neighborhood in Lemma
\ref{lem:endpoint-point-saddle}.  Uniformly over
$\mathfrak Q_N(\ell,\eta,M)$ and $h\in\mathcal U_n$,
\[
        cD_q\le\Sigma(h)\le CD_q
\tag{10.11}\label{eq:endpoint-point-covariance-comparison}
\]
in quadratic-form order.  For every fixed $A<\infty$, uniformly for lattice
points $v\in\mathbb Z^\ell$ satisfying
\[
        \|D_q^{-1/2}(v-M(h))\|\le A,
\]
we have
\[
\begin{aligned}
\Pp(Z(h)=v)
&=(2\pi)^{-\ell/2}(\det\Sigma(h))^{-1/2}\\
&\quad\times
\exp\left[-\frac12(v-M(h))^T\Sigma(h)^{-1}(v-M(h))\right]
+o\left(\prod_{j=1}^\ell q_j^{-1/2}\right).
\end{aligned}
\tag{10.12}\label{eq:endpoint-point-llt}
\]
For every fixed $\nu<1/6$, there are $C,c>0$ such that
\[
\Pp(Z(h)=v)
\le C\prod_{j=1}^\ell q_j^{-1/2}
\exp\left[-c\sum_{j=1}^\ell
\frac{(v_j-M_j(h_j))^2}{q_j}\right]
\tag{10.13}\label{eq:endpoint-point-local-upper}
\]
whenever
$|v_j-M_j(h_j)|\le q_j^{1/2+\nu}$ for every $j$.
\end{lemma}

\begin{proof}
For $b\notin P$, the vector
\[
        V_b(h)=
        (\one_{\{Y_{n,b}<y_1(h_1)\}},\ldots,
         \one_{\{Y_{n,b}<y_\ell(h_\ell)\}})
\]
has the cumulative categorical support
\[
        (1,\ldots,1),(0,1,\ldots,1),\ldots,
        (0,\ldots,0,1),0.
\]
By \eqref{eq:endpoint-point-chamber-margin}, the adjacent threshold gaps are
uniformly positive in $\mathcal U_n$.  Choose a small multiplicative interval
$[(1-\delta)q_j,(1+\delta)q_j]$ for each $j$, with $\delta$ so small that the
$\ell$ intervals are disjoint.  Uniform regular variation and
\eqref{eq:endpoint-point-saddle-asymptotics} imply that, for every index $b$ in
the $j$-th block, the two categories immediately below and above the threshold
$y_j(h_j)$ both have probability bounded below by a positive constant.
Removing the fixed set $P$ changes each block size by at most $\ell$.
Consequently, for every $j$, there are at least $cq_j$ independent summands
having two atoms whose difference is the coordinate vector $e_j$, with both
atom probabilities at least $c$.  These bounds persist under every bounded
exponential tilt.

The disjoint blocks give the covariance lower bound in
\eqref{eq:endpoint-point-covariance-comparison}.  The upper bound follows from
\[
        \Var Z_j(h)=O(q_j),
\]
which is Lemma \ref{lem:endpoint-rv-sums} applied at the relevant compact
values of $e^{-h_j}$, followed by Cauchy's inequality.  The same blocks and
Lemma \ref{lem:coord-fourier} give, uniformly under bounded tilts,
\[
\left|\Ee e^{iu\cdot(Z(h)-M(h))}\right|
\le\exp\left[-c\sum_{j=1}^\ell
q_j\operatorname{dist}(u_j,2\pi\mathbb Z)^2\right].
\tag{10.14}\label{eq:endpoint-point-fourier}
\]
On the anisotropic box $|u_j|\le A q_j^{-1/2}$, bounded categorical
increments and the covariance upper bound give
\[
\log\Ee e^{iu\cdot(Z-M)}
=-\frac12u^T\Sigma(h)u
+O\left((\max_j|u_j|)\sum_jq_ju_j^2\right).
\]
After the change of variables $u=D_q^{-1/2}w$, the remainder is
$O_A(q_1^{-1/2})=o(1)$.  Fourier inversion and
\eqref{eq:endpoint-point-fourier} control the complementary region on the
product scale.  This proves \eqref{eq:endpoint-point-llt}.  The tilted local
upper bound \eqref{eq:endpoint-point-local-upper} follows from Lemma
\ref{lem:anisotropic-tilted-upper}: the tilted covariance upper bound follows
from bounded distortion of the categorical probabilities, and the edge blocks
supply the required tilted concentration estimate.
\end{proof}

\begin{lemma}[Source-to-target mixed localization]
\label{lem:endpoint-point-localization}
There is $L_0=L_0(\ell,\eta,M,\gamma)$ such that the following holds for
$L\ge L_0$.  For $h\in\mathscr C_{n,q}$, set
\[
\begin{aligned}
        B_R(h)&=\{j:h_j>L\},\\
        B_L(h)&=\{j:h_j<-L\},\\
        G(h)&=\{1,\ldots,\ell\}\setminus(B_R(h)\cup B_L(h)),
\end{aligned}
\]
and put
\[
        H_j=z_j-M_j(h_j),\qquad H_0=0.
\]
On every region on which these three sets are fixed,
\[
\begin{aligned}
\Pp(Z(h)=z)
&\le C_L\prod_{j=1}^\ell q_j^{-1/2}
\exp\left[-c\sum_{j\in B_R}|H_j|
           -c\sum_{j\in B_L}q_j\right]\\
&\quad\times
\exp\left[-c_L\sum_{j\in G}
\min\left\{\frac{H_j^2}{q_j},q_j\right\}\right].
\end{aligned}
\tag{10.15}\label{eq:endpoint-point-mixed-bound}
\]
Consequently,
\[
\begin{aligned}
\Pp(Z(h)=z)
&\le C_L\prod_{j=1}^\ell q_j^{-1/2}
\exp\left[-ce^{L/(2\gamma)}\sum_{j\in B_R}q_j
           -c\sum_{j\in B_L}q_j\right]\\
&\quad\times
\exp\left[-c_L\sum_{j\in G}
q_j\min\{(h_j-h_j^*)^2,1\}\right].
\end{aligned}
\tag{10.16}\label{eq:endpoint-point-scaled-bound}
\]
Moreover,
\[
\lim_{L\to\infty}\limsup_{N\to\infty}
\sup_{\mathfrak Q_N(\ell,\eta,M)}
\frac{
\displaystyle
\int_{\mathscr C_{n,q}\cap\{h\notin[-L,L]^\ell\}}
\prod_{j=1}^\ell\varphi_{R_{n,j}}(h_j)
\Pp(Z(h)=z)\dd h
}{\prod_{j=1}^\ell q_j^{-1}}=0.
\tag{10.17}\label{eq:endpoint-point-noncompact}
\]
On every fixed compact cube containing the saddle vectors,
\[
\Pp(Z(h)=z)
\le C\prod_{j=1}^\ell q_j^{-1/2}
\exp\left[-c\sum_{j=1}^\ell
q_j\min\{(h_j-h_j^*)^2,1\}\right].
\tag{10.18}\label{eq:endpoint-point-compact-localization}
\]
\end{lemma}

\begin{proof}
Define interval counts
\[
        W_1=Z_1,\qquad W_j=Z_j-Z_{j-1}\quad(2\le j\le\ell),
\]
with targets
\[
        w_1=q_1-1,\qquad
        w_j=q_j-q_{j-1}-1\quad(2\le j\le\ell).
\]
The target separation gives $c_\eta q_j\le w_j\le q_j$.  Let
$m_j=\Ee W_j$ and $d_j=w_j-m_j$; then $d_j=H_j-H_{j-1}$.
For $j\in G$, one has $m_j\le M_j(h_j)=O_L(q_j)$, and hence
$|d_j|=O_L(q_j)$.

Fix the offender sets.  Let
\[
        K_h(\theta)=\log\Ee
        \exp\{\theta\cdot(W(h)-m(h))\}
\]
be the centered interval-coordinate cumulant, and let
\[
        (\mathcal T\alpha)_k=\sum_{j\ge k}\alpha_j.
\]
Then
$(\mathcal T\alpha)\cdot d=\alpha\cdot H$.
Choose $\delta>0$ sufficiently small and set
\[
        \alpha_j^-=-\delta\one_{\{j\in B_R\}},
        \qquad
        \alpha_j^+=\delta\one_{\{j\in B_L\}}.
\]
On $G$, use the local interval tilt
\[
        \lambda_j=
        \begin{cases}
        \delta_Ld_j/q_j,& |d_j|\le\varepsilon q_j,\\
        \delta_L\operatorname{sgn}(d_j),& |d_j|>\varepsilon q_j,
        \end{cases}
\]
where $\varepsilon<c_\eta/4$ and $\delta_L>0$ is sufficiently small; put
$\lambda_j=0$ off $G$.  Use the single total interval-coordinate tilt
\[
        \Theta=\mathcal T\alpha^-+\mathcal T\alpha^++\lambda.
\tag{10.19}\label{eq:endpoint-point-total-tilt}
\]
Generalized H\"older inequality yields
\[
\begin{aligned}
K_h(\Theta)
&\le\frac13K_h(3\mathcal T\alpha^-)
 +\frac13K_h(3\mathcal T\alpha^+)
 +\frac13K_h(3\lambda).
\end{aligned}
\tag{10.20}\label{eq:endpoint-point-holder}
\]

By \eqref{eq:endpoint-point-mean-asymptotic}, for large $L$,
\[
\begin{array}{lll}
 j\in B_R:&M_j(h_j)\ge ce^{L/\gamma}q_j,&H_j<0,\\[1mm]
 j\in B_L:&M_j(h_j)\le Ce^{-L/\gamma}q_j,&0<H_j\le q_j.
\end{array}
\tag{10.21}\label{eq:endpoint-point-offender-geometry}
\]
The proof of Lemma
\ref{lem:actual-mean-lower-cumulative-chernoff} applies to the present nested
lower-priority counts without change.  Applied with parameter $3\delta$, it
gives
\[
        -3\alpha^-\cdot H+K_h(3\mathcal T\alpha^-)
        \le-c\sum_{j\in B_R}|H_j|,
\tag{10.22}\label{eq:endpoint-point-right-cost}
\]
\[
        -3\alpha^+\cdot H+K_h(3\mathcal T\alpha^+)
        \le-c\sum_{j\in B_L}q_j.
\tag{10.23}\label{eq:endpoint-point-left-cost}
\]
Under every bounded interval tilt, the weighted-sum estimate gives
\[
        \Var_\xi(v\cdot W(h))\le C_L\sum_{j\in G}q_jv_j^2
\]
for $v$ supported on $G$.  Taylor's formula for the cumulant and the choice of
$\delta_L$ therefore give
\[
        -3\lambda\cdot d+K_h(3\lambda)
        \le-c_L\sum_{j\in G}
        \min\left\{\frac{d_j^2}{q_j},q_j\right\}.
\tag{10.24}\label{eq:endpoint-point-local-cost}
\]

All coordinates of $\Theta$ are bounded for fixed $L$.  For every locally
regular coordinate $j\in G$ with $|d_j|\le\varepsilon q_j$, the interval mean
is at least $(c_\eta/2)q_j$.  The middle-block argument in Lemma
\ref{lem:endpoint-interval-edge-blocks}, with $q_j$ in place of $a_j$ and with
the fixed deletion set $P$, supplies forward two-atom blocks of order $q_j$.
They can be chosen disjointly, and persist under the total tilt.  After
projection onto the locally regular coordinates, Lemma
\ref{lem:projected-forward-edge-concentration} gives
\[
        \Pp_\Theta(W_{G_0}=w_{G_0})
        \le C_L\prod_{j\in G_0}q_j^{-1/2}.
\tag{10.25}\label{eq:endpoint-point-tilted-factor}
\]
The exact change-of-measure identity, together with
\eqref{eq:endpoint-point-holder}--\eqref{eq:endpoint-point-tilted-factor},
gives the desired exponential estimate first in terms of the interval
residuals $d_j$.  Indices lacking a local factor carry a fixed exponential
loss of order $q_j$, which absorbs the missing factor $q_j^{-1/2}$.

To pass from interval to cumulative residuals without losing boundary terms,
apply Lemma \ref{lem:boundary-truncated-chain} with $A_j=q_j$, with its
$B_-$ equal to the present $B_R$, and its $B_+$ equal to the present $B_L$.
The assumptions follow from $|d_j|=O_L(q_j)$ on $G$ and
$|H_j|\le q_j$ on $B_L$.  Reserving a fixed fraction of the offender costs
proves \eqref{eq:endpoint-point-mixed-bound}.  On $G$, the derivative bounds
and monotonicity imply
\[
        |H_j|\ge c_Lq_j\min\{|h_j-h_j^*|,1\}.
\]
On $B_R$, \eqref{eq:endpoint-point-offender-geometry} gives
$|H_j|\ge ce^{L/\gamma}q_j$.  This proves
\eqref{eq:endpoint-point-scaled-bound}; the compact estimate
\eqref{eq:endpoint-point-compact-localization} is the case with no offenders.

It remains to integrate the tail.  By
\eqref{eq:endpoint-point-rate-ratio}, $R_{n,j}$ lies in a fixed compact
subinterval of $(0,\infty)$.  Uniformly in such $R$,
\[
        \varphi_R(h)\le Ce^{-h}\quad(h\ge0),
        \qquad
        \varphi_R(h)\le Ce^{-h}e^{-ce^{-h}}\quad(h\le0).
\tag{10.26}\label{eq:endpoint-point-density-tails}
\]
Each nonoffending coordinate contributes $O(q_j^{-1/2})$ after integration on
its Gaussian scale, in addition to the prefactor $q_j^{-1/2}$.  For every
offender, the exponential loss in \eqref{eq:endpoint-point-scaled-bound}
absorbs the missing local factor, and the remaining density tail tends to zero
as $L\to\infty$.  Summing over the finitely many offender sets proves
\eqref{eq:endpoint-point-noncompact}.
\end{proof}

\begin{lemma}[Source-to-target covariance stability]
\label{lem:endpoint-point-covariance-stability}
Let
\[
        H_q=\diag(q_1^{-1/2},\ldots,q_\ell^{-1/2}),
        \qquad \Sigma_*=\Sigma(h^*).
\]
For every fixed $A<\infty$,
\[
\sup_{\|u\|\le A}
\left\|
D_q^{-1/2}
\bigl[\Sigma(h^*+H_qu)-\Sigma_*\bigr]
D_q^{-1/2}
\right\|_{\mathrm{op}}\longrightarrow0
\tag{10.27}\label{eq:endpoint-point-covariance-stability}
\]
uniformly over $\mathfrak Q_N(\ell,\eta,M)$.  Consequently,
\[
        \frac{\det\Sigma(h^*+H_qu)}{\det\Sigma_*}\longrightarrow1
\tag{10.28}\label{eq:endpoint-point-determinant-stability}
\]
and
\[
\left\|
D_q^{1/2}
\bigl[\Sigma(h^*+H_qu)^{-1}-\Sigma_*^{-1}\bigr]
D_q^{1/2}
\right\|_{\mathrm{op}}\longrightarrow0
\tag{10.29}\label{eq:endpoint-point-inverse-stability}
\]
uniformly on the same sets.
\end{lemma}

\begin{proof}
For $j\le k$, put
\[
        p_{b,j}(h_j)=
        \exp\left(-e^{-h_j}
        \frac{\lambda_{n,b}}{\lambda_{n,q_j}}\right).
\]
The threshold events are nested in the chamber, so
\[
        \Sigma_{jk}(h)
        =\sum_{b\notin P}p_{b,j}(h_j)(1-p_{b,k}(h_k)).
\tag{10.30}\label{eq:endpoint-point-covariance-entry}
\]
Lemma \ref{lem:endpoint-rv-sums}, with power one, gives on a fixed saddle
neighborhood
\[
        |\partial_{h_j}\Sigma_{jk}(h)|\le Cq_j,
        \qquad
        |\partial_{h_k}\Sigma_{jk}(h)|\le Cq_k.
\]
Hence, for $\|u\|\le A$,
\[
        |\Sigma_{jk}(h^*+H_qu)-\Sigma_{jk}(h^*)|
        \le C_A(\sqrt{q_j}+\sqrt{q_k}).
\]
After division by $\sqrt{q_jq_k}$ this is at most
$C_A(q_j^{-1/2}+q_k^{-1/2})\le2C_Aq_1^{-1/2}$.  Since $\ell$ is fixed,
\eqref{eq:endpoint-point-covariance-stability} follows.  The normalized
covariance matrices remain in a compact subset of the positive-definite
matrices by \eqref{eq:endpoint-point-covariance-comparison}; continuity of
determinant and inversion gives
\eqref{eq:endpoint-point-determinant-stability}--
\eqref{eq:endpoint-point-inverse-stability}.
\end{proof}

\begin{lemma}[Endpoint integrated point-constraint saddle]
\label{lem:endpoint-general-integrated-saddle}
Uniformly over $\mathfrak Q_N(\ell,\eta,M)$,
\[
\begin{aligned}
&\int_{\mathscr C_{n,q}}
\prod_{j=1}^\ell\varphi_{R_{n,j}}(h_j)
\Pp(Z(h)=z)\dd h\\
&\qquad=(1+o(1))
\prod_{j=1}^\ell
\frac{\varphi_{R_{n,j}}(h_j^*)}{d_{n,j}},
\end{aligned}
\tag{10.31}\label{eq:endpoint-point-integrated-saddle}
\]
where $d_{n,j}=M_j'(h_j^*)$.
\end{lemma}

\begin{proof}
By Lemma \ref{lem:endpoint-point-localization}, it suffices first to integrate
over a fixed compact cube and then over a shrinking saddle neighborhood; the
complement is $o(\prod_jq_j^{-1})$.  Put
\[
        h=h^*+H_qu,
        \qquad
        D_n=\diag(d_{n,1},\ldots,d_{n,\ell}).
\]
For bounded $u$, Taylor's theorem and
\eqref{eq:endpoint-point-mean-derivatives} give
\[
        M(h)-z=D_nH_qu+o(D_q^{1/2})
\tag{10.32}\label{eq:endpoint-point-mean-taylor}
\]
uniformly.  The density factors satisfy
\[
        \prod_j\varphi_{R_{n,j}}(h_j)
        =\prod_j\varphi_{R_{n,j}}(h_j^*)(1+o(1))
\]
on bounded $u$-balls.  Lemma \ref{lem:endpoint-point-llt} and Lemma
\ref{lem:endpoint-point-covariance-stability} therefore give, uniformly on
such balls, the Gaussian integrand with covariance $\Sigma_*$ and mean shift
$D_nH_qu$.

Since $\dd h=|\det H_q|\dd u$, the Gaussian integral contributes
\[
\begin{aligned}
&(2\pi)^{-\ell/2}(\det\Sigma_*)^{-1/2}|\det H_q|\\
&\quad\times
\int_{\mathbb R^\ell}
\exp\left[-\frac12u^TH_q^TD_n^T
\Sigma_*^{-1}D_nH_qu\right]\dd u
=|\det D_n|^{-1}.
\end{aligned}
\tag{10.33}\label{eq:endpoint-point-determinant-cancellation}
\]
Indeed,
\[
\det(H_q^TD_n^T\Sigma_*^{-1}D_nH_q)
=|\det H_q|^2|\det D_n|^2(\det\Sigma_*)^{-1}.
\]
The compact localization bound
\eqref{eq:endpoint-point-compact-localization} is an integrable Gaussian
majorant in $u$, and
\eqref{eq:endpoint-point-noncompact} handles the outer region.  Dominated
convergence proves \eqref{eq:endpoint-point-integrated-saddle}.
\end{proof}

\begin{proposition}[Separated endpoint entry estimate]
\label{prop:endpoint-general-entry}
Uniformly over $\mathfrak Q_N(\ell,\eta,M)$,
\[
        \Pp(\tau_n(p_j)=q_j,\ 1\le j\le\ell)
        =(1+o(1))\prod_{j=1}^\ell
        \frac1{q_j}k_\gamma\!\left(\log\frac{q_j}{p_j}\right).
\tag{10.34}\label{eq:endpoint-general-entry}
\]
\end{proposition}

\begin{proof}
Condition on the specified centered Gumbels $G_{p_j}$ and put
\[
        h_j=\log R_{n,j}+G_{p_j}.
\]
Then the specified priority is
$Y_{n,p_j}=\log\lambda_{n,q_j}+h_j$, and $h_j$ has density
$\varphi_{R_{n,j}}$.  Since $q_1<\cdots<q_\ell$, the point constraints force
$Y_{n,p_1}<\cdots<Y_{n,p_\ell}$.  In this chamber exactly $j-1$ specified
priorities lie below the $j$-th threshold, so the event is equivalent to
$Z_j(h)=q_j-j$ for all $j$.  Thus the probability is exactly the integral in
\eqref{eq:endpoint-point-integrated-saddle}.

By Lemma \ref{lem:endpoint-point-saddle},
$h_j^*\to-\log K_\gamma$ and $d_{n,j}\sim q_j/\gamma$.  Therefore
\[
\begin{aligned}
\frac{\varphi_{R_{n,j}}(h_j^*)}{d_{n,j}}
&=\frac{\gamma K_\gamma}{q_j}
R_{n,j}e^{-K_\gamma R_{n,j}}(1+o(1))\\
&=\frac1{q_j}k_\gamma\!\left(\log\frac{q_j}{p_j}\right)(1+o(1)),
\end{aligned}
\]
where the second equality uses
\eqref{eq:endpoint-point-rate-ratio}.  Multiply over $j$.
\end{proof}

\subsection{A global cycle envelope}

The next lemma makes precise the monotonicity used when forced deletions are
interspersed with ordinary Luce deletions.

\begin{lemma}[Scheduled-deletion survival]
\label{lem:scheduled-deletion-survival}
Let $S$ be a finite set of labels with positive rates, and let $p\in S$.  Start
from $S$ and run a schedule in which each step is either an ordinary Luce
deletion from the current set or a forced removal of a chosen competitor
$c\ne p$.  The forced choices may depend on the previously exposed history.
For $m\ge0$, the probability that $p$ is still present after $m$ ordinary
deletions is at most
\[
        H_m(S,p):=
        \Pp\left(\#\{b\in S\setminus\{p\}:T_b<T_p\}\ge m\right)
\tag{10.35}\label{eq:scheduled-survival-function}
\]
for an unmodified exponential race on $S$.
\end{lemma}

\begin{proof}
The function $H_m(S,p)$ is increasing under addition of competitors, by
coupling the clocks.  It also satisfies the exact recursion
\[
        H_m(S,p)=
        \sum_{b\in S\setminus\{p\}}
        \frac{\lambda_b}{\sum_{c\in S}\lambda_c}
        H_{m-1}(S\setminus\{b\},p),
        \qquad m\ge1,
\tag{10.36}\label{eq:scheduled-survival-recursion}
\]
with $H_0=1$: to survive $m$ ordinary deletions, the first ordinary deletion
must be a competitor $b$, after which memorylessness leaves the same problem
with $m-1$ deletions.

Induct on the remaining schedule.  If the next step is a forced removal of
$c$, the induction hypothesis and monotonicity give a bound by
$H_m(S\setminus\{c\},p)\le H_m(S,p)$.  If the next step is ordinary, condition
on the deleted label and use \eqref{eq:scheduled-survival-recursion}.  This
also covers adaptive forced choices after conditioning on the current history.
\end{proof}

\begin{lemma}[Forced-deletion bound]
\label{lem:forced-deletion}
Consider compatible point constraints
$\tau_n(p_i)=q_i$, $1\le i\le L$, with distinct sources and distinct targets.
Then
\[
        \Pp(\tau_n(p_i)=q_i,\ 1\le i\le L)
        \le C^L\prod_{i=1}^L
        \frac{\lambda_{n,p_i}}{q_i\lambda_{n,q_i}}.
\tag{10.37}\label{eq:forced-deletion-basic}
\]
If $(p_*,q_*)$ is one distinguished constraint, the right-hand side may be
multiplied by
\[
        \Pp(\tau_n(p_*)\le q_*+L).
\tag{10.38}\label{eq:forced-deletion-survival-factor}
\]
\end{lemma}

\begin{proof}
Expose the increasing exponential race as a pure-death chain.  When $q$ labels
are alive, every alive set has total rate at least
\[
        \sum_{b\le q}\lambda_{n,b}
        \ge c_\Sigma q\lambda_{n,q}
\]
by Lemma \ref{lem:endpoint-partial-sum}.  Hence the conditional probability of
deleting source $p_i$ at level $q_i$ is at most
$C\lambda_{n,p_i}/(q_i\lambda_{n,q_i})$.  Sum over compatible death histories
and pull out these bounds at the forced levels.  The remaining transition
weights have total mass at most one, proving
\eqref{eq:forced-deletion-basic}.

For the refinement, retain the condition that $p_*$ survives until level
$q_*$.  Before that level at most $L-1$ deletions are forced.  Thus the
residual schedule contains at least $n-q_*-L$ ordinary deletions before
$p_*$ may be removed.  Lemma \ref{lem:scheduled-deletion-survival} bounds the
residual path mass by
\[
\begin{aligned}
H_{n-q_*-L}([n],p_*)
&=\Pp\left(\#\{b\ne p_*:T_b<T_{p_*}\}
\ge n-q_*-L\right)\\
&=\Pp(\tau_n(p_*)\le q_*+L),
\end{aligned}
\]
which proves \eqref{eq:forced-deletion-survival-factor}.
\end{proof}

\begin{lemma}[Endpoint survival tail]
\label{lem:endpoint-survival-tail}
Put
\[
        U_n=\left\lfloor\frac{n}{(\log n)^2}\right\rfloor,
        \qquad V_n=2U_n.
\]
There are constants $C,c,\delta>0$ such that, uniformly for
$1\le q<p\le n$, $q\le U_n$, and every fixed $L$,
\[
        \Pp(\tau_n(p)\le q+L)
        \le C\exp\left\{-c
        \left(\frac{p\wedge V_n}{q+L}\right)^\delta\right\}.
\tag{10.39}\label{eq:endpoint-survival-tail}
\]
\end{lemma}

\begin{proof}
Put $m=q+L$, $\widetilde p=p\wedge V_n$, and choose
$\beta\in(0,\gamma)$ and $\delta=\beta/(1+\beta)$.  If
$d=\widetilde p/m$ is bounded, enlarge $C$.  Otherwise let
$s=\lfloor m d^\delta\rfloor$.  Then $Cm\le s\le\widetilde p/2$, so in
particular the block $\{1,\ldots,s\}$ does not contain $p$.  Take
$t=1/\lambda_{n,s}$.  For every $b\le s$,
\[
        \Pp(T_{n,b}>t)=e^{-\lambda_{n,b}/\lambda_{n,s}}\ge e^{-1}.
\]
If $T_{n,p}<t$ and at least $m$ of these first $s$ clocks exceed $t$, then at
least $m$ clocks exceed $T_{n,p}$, so $\tau_n(p)\ge m+1$.  Therefore
\[
        \Pp(\tau_n(p)\le m)
        \le e^{-\lambda_{n,p}/\lambda_{n,s}}+e^{-cs}
\]
by a binomial Chernoff bound.  Potter's lower bound and monotonicity give
\[
        \lambda_{n,p}/\lambda_{n,s}
        \ge c(\widetilde p/s)^\beta.
\]
By the choice of $s$, both $s$ and $(\widetilde p/s)^\beta$ are bounded below
by a constant times $d^\delta$.  This proves
\eqref{eq:endpoint-survival-tail}.
\end{proof}

Let $\Gamma$ be a prescribed collection of pairwise vertex-disjoint directed
cycles in terminal coordinates, and let $E_\Gamma$ be the event that all those
cycles occur.  For a cycle component $\gamma$ define its capped logarithmic
oscillation by
\[
        \omega_{U_n}(\gamma)
        =\operatorname{osc}_{a\in\gamma}\log(a\wedge U_n).
\tag{10.40}\label{eq:capped-log-oscillation}
\]

\begin{lemma}[Capped path drop]
\label{lem:capped-path-drop}
Let $\gamma$ be a nontrivial directed cycle containing at least one vertex not
exceeding $U_n$.  There is a downward edge $p\mapsto q$ of $\gamma$ such that
$q\le U_n$ and
\[
        \frac{p\wedge V_n}{q}
        \ge\exp\{\omega_{U_n}(\gamma)/|\gamma|\}.
\tag{10.41}\label{eq:capped-path-drop}
\]
\end{lemma}

\begin{proof}
If $\omega_{U_n}(\gamma)=0$, then the component must contain a vertex above
$U_n$ and a vertex equal to $U_n$; any directed entrance from the upper part
to a vertex not exceeding $U_n$ gives a downward edge satisfying
\eqref{eq:capped-path-drop}.  Assume henceforth that the oscillation is
positive.  Choose vertices whose capped logarithms attain the maximum and
minimum, and follow the directed cycle from the former to the latter.  Put
$z_i=v_i\wedge U_n$ along this path.  The product of the ratios
$z_i/z_{i+1}$ over the downward capped steps is at least the net ratio between
the endpoints, namely $e^{\omega_{U_n}(\gamma)}$.  There are at most
$|\gamma|$ steps, so one downward capped step has ratio at least the right-hand
side of \eqref{eq:capped-path-drop}.  For that step,
$z_{i+1}=q\le U_n$ and
$p\wedge V_n\ge p\wedge U_n=z_i$, which proves the claim.
\end{proof}

\begin{corollary}[Cycle envelope]
\label{cor:endpoint-cycle-envelope}
If $L$ is the total number of vertices of $\Gamma$, then
\[
        \Pp(E_\Gamma)
        \le\frac{C_L}{\prod_{a\in V(\Gamma)}a}.
\tag{10.42}\label{eq:cycle-envelope-basic}
\]
If $\gamma$ is a nontrivial component meeting $[1,U_n]$, then
\[
        \Pp(E_\Gamma)
        \le\frac{C_L}{\prod_{a\in V(\Gamma)}a}
        \Psi_L(\omega_{U_n}(\gamma)),
\tag{10.43}\label{eq:cycle-envelope-component}
\]
where
\[
        \Psi_L(x)=C_L\exp\{-c_Le^{c_Lx}\}.
\tag{10.44}\label{eq:cycle-envelope-psi}
\]
Consequently, if
\[
        \Omega_{U_n}(\Gamma)
        =\max_{\substack{\gamma\text{ nontrivial}\
                         \gamma\cap[1,U_n]\ne\varnothing}}
        \omega_{U_n}(\gamma),
\]
then, whenever the maximum is over a nonempty set,
\[
        \Pp(E_\Gamma)
        \le\frac{C_L}{\prod_{a\in V(\Gamma)}a}
        \Psi_L(\Omega_{U_n}(\Gamma)).
\tag{10.45}\label{eq:cycle-envelope-maximum}
\]
\end{corollary}

\begin{proof}
Apply Lemma \ref{lem:forced-deletion}.  Around every directed cycle the rate
ratios telescope:
\[
        \prod_{a\mapsto b}
        \frac{\lambda_{n,a}}{b\lambda_{n,b}}
        =\frac1{\prod_{a\in\gamma}a}.
\]
This proves \eqref{eq:cycle-envelope-basic}.  For a component meeting
$[1,U_n]$, choose the edge supplied by Lemma \ref{lem:capped-path-drop} as the
distinguished constraint in Lemma \ref{lem:forced-deletion}.  Lemma
\ref{lem:endpoint-survival-tail} gives
\eqref{eq:cycle-envelope-component}.  Choosing a component attaining the
maximum gives \eqref{eq:cycle-envelope-maximum}.
\end{proof}

\subsection{The logarithmic trace}

\begin{lemma}[Continuum multiplicative trace]
\label{lem:continuum-endpoint-trace}
For every fixed $r\ge1$,
\[
        I_r(T):=\int_{[0,T]^r}
        \prod_{j=1}^r k_\gamma(x_{j+1}-x_j)
        \dd x_1\cdots\dd x_r
        =T k_\gamma^{*r}(0)+O_r(1),
\tag{10.46}\label{eq:continuum-endpoint-trace}
\]
where $x_{r+1}=x_1$.
\end{lemma}

\begin{proof}
For $r=1$ the identity is exact.  For $r\ge2$, put
$z_j=x_{j+1}-x_j$, $1\le j<r$, and
$z_r=-\sum_{j<r}z_j$.  If $s_0=0$ and
$s_j=z_1+\cdots+z_j$, let
$\mathcal R(z)=\max_js_j-\min_js_j$.  The admissible translations $x_1$ have
length $[T-\mathcal R(z)]_+$, so
\[
        I_r(T)=\int_{\mathbb R^{r-1}}[T-\mathcal R(z)]_+
        \prod_{j=1}^r k_\gamma(z_j)\dd z_1\cdots\dd z_{r-1}.
\]
The product integral without the bracket is $k_\gamma^{*r}(0)$.  The density
$k_\gamma$ has an exponential right tail, a superexponential left tail, and a
finite first absolute moment.  Since
$\mathcal R(z)\le2\sum_{j<r}|z_j|$, the error integral is finite uniformly in
$T$.
\end{proof}

We shall use the elementary harmonic-coordinate bounds
\[
        \sum_{A\le a\le B}\frac1a
        \le1+\log(B/A),
\tag{10.47}\label{eq:harmonic-interval-bound}
\]
and, for $J_a=[\log a,\log(a+1))$,
\[
        |J_a|=\frac1a+O(a^{-2}).
\tag{10.48}\label{eq:log-cell-length}
\]

\begin{lemma}[Discrete endpoint trace]
\label{lem:discrete-endpoint-trace}
Let $L_n\to\infty$, $U_n=o(n)$, and
$T_n=\log(U_n/L_n)\to\infty$.  Then
\[
\begin{aligned}
&\sum_{\substack{L_n\le a_1,\ldots,a_r\le U_n\\
        a_1,\ldots,a_r\text{ distinct}}}
        \prod_{j=1}^r\frac1{a_{j+1}}
        k_\gamma\!\left(\log\frac{a_{j+1}}{a_j}\right)\\
&\qquad=k_\gamma^{*r}(0)T_n+o(T_n),
\end{aligned}
\tag{10.49}\label{eq:discrete-endpoint-trace}
\]
where $a_{r+1}=a_1$.
\end{lemma}

\begin{proof}
Because $\prod_ja_{j+1}^{-1}=\prod_ja_j^{-1}$, the sum is a logarithmic
Riemann sum.  First restrict to
$\operatorname{osc}(\log a_1,\ldots,\log a_r)\le M$.  On this region the
cyclic kernel product and its first derivatives are bounded by a constant
depending on $r,M$.  Using the cells $J_a$ in
\eqref{eq:log-cell-length}, cell integration and
\eqref{eq:harmonic-interval-bound} show that the difference from the continuum
integral over $[\log L_n,\log U_n]^r$ is
\[
        O_{r,M}(T_n/L_n)=o(T_n).
\tag{10.50}\label{eq:trace-riemann-error}
\]
The same bound controls the replacement of $1/a$ by $|J_a|$.

We next record a uniform discrete tail bound.  The local supremum of
$k_\gamma$ over an interval of length one is bounded by an integrable density
$\overline k_\gamma$ having the same exponential right tail, a
superexponential left tail, and a finite first moment.  Enlarging every
logarithmic cell by one and using
\eqref{eq:harmonic-interval-bound}, the discrete cyclic sum with oscillation
larger than $M$ is bounded by a constant times the corresponding continuum
cyclic integral with $k_\gamma$ replaced by $\overline k_\gamma$ and
oscillation larger than $M-r$.  The increment representation in Lemma
\ref{lem:continuum-endpoint-trace} therefore gives
\[
\lim_{M\to\infty}\sup_n\frac1{T_n}
\sum_{\substack{L_n\le a_1,\ldots,a_r\le U_n\\
 \operatorname{osc}(\log a)>M}}
\prod_{j=1}^r\frac1{a_{j+1}}
 k_\gamma\!\left(\log\frac{a_{j+1}}{a_j}\right)=0.
\tag{10.51}\label{eq:trace-discrete-tail}
\]

Finally, an equality $a_i=a_j$ replaces one harmonic factor by a square:
$\sum_{a\ge L_n}a^{-2}=O(L_n^{-1})$.  The remaining localized relative sums
are uniformly bounded by the same convolution majorant, so all equality
diagonals contribute $O_r(L_n^{-1})$.  Combine
\eqref{eq:trace-riemann-error}--\eqref{eq:trace-discrete-tail} with Lemma
\ref{lem:continuum-endpoint-trace}, and then let $M\to\infty$.
\end{proof}

\subsection{Factorial-moment summation}

The next elementary observation is used repeatedly.  If a component has
$r$ vertices, one may choose one logarithmic translation coordinate and
$r-1$ relative coordinates.  By \eqref{eq:harmonic-interval-bound}, if its
logarithmic oscillation is at most $x$, the total harmonic mass of the relative
coordinates is at most $C_r(1+x)^{r-1}$.  For an ordered collection of $m$
components and $L$ vertices, the corresponding bound is
\[
        C_L(1+x)^{L-m}(1+T_n)^m.
\tag{10.52}\label{eq:component-harmonic-volume}
\]
If one component is required to meet $[1,L_n)$, one of the $T_n$ translation
factors in \eqref{eq:component-harmonic-volume} is replaced by
$O(1+\log L_n+x)$.  If a component crosses $U_n$ and its capped oscillation is
at most $x$, its lower anchor belongs to $[U_ne^{-x},U_n]$, so its translation
factor is $O(1+x)$ rather than $T_n$; each vertex above $U_n$ contributes the
harmonic factor
\[
        H_n:=\sum_{a>U_n}\frac1a=O(\log\log n).
\tag{10.53}\label{eq:upper-harmonic-mass}
\]
These bounds follow by summing the nonanchor vertices successively over
multiplicative intervals and will be used without further comment.

\begin{lemma}[Negligible endpoint cycle configurations]
\label{lem:negligible-endpoint-cycle-configurations}
Fix an ordered collection of $m$ cycle components with total number of vertices
$L$.  Put
\[
        L_n=\lceil\log n\rceil,
        \qquad
        U_n=\left\lfloor\frac{n}{(\log n)^2}\right\rfloor,
        \qquad
        T_n=\log(U_n/L_n).
\]
Then:
\begin{enumerate}[label=(\roman*)]
\item among collections with all vertices in $[L_n,U_n]$, the total probability
of those for which some nontrivial component has logarithmic oscillation greater
than $M$ is at most $\varepsilon_M T_n^m$, where
$\varepsilon_M\downarrow0$;
\item if all component oscillations are at most $M$, the total probability of
collections containing two vertices at logarithmic distance at most $\eta$ is
at most
\[
        C_{L,M}\eta T_n^m+O_{L,M}(T_n^{m-1});
\tag{10.54}\label{eq:close-cycle-configurations}
\]
\item collections having at least one vertex outside $[L_n,U_n]$ contribute
$o(T_n^m)$.
\end{enumerate}
\end{lemma}

\begin{proof}
For (i), let $\Omega$ be the maximum logarithmic oscillation among the
nontrivial components.  Since all vertices are at most $U_n$, this is the same
as $\Omega_{U_n}(\Gamma)$.  Corollary
\ref{cor:endpoint-cycle-envelope} and
\eqref{eq:component-harmonic-volume} give, after decomposing
$\Omega\in[k,k+1)$,
\[
\begin{aligned}
&\sum_{\substack{\Gamma:\ V(\Gamma)\subset[L_n,U_n]\\
                  \Omega(\Gamma)>M}}
\Pp(E_\Gamma)\\
&\qquad\le C_LT_n^m
\sum_{k\ge M}(2+k)^{L-m}\Psi_L(k)
=:\varepsilon_MT_n^m.
\end{aligned}
\]
The superexponential decay in \eqref{eq:cycle-envelope-psi} gives
$\varepsilon_M\downarrow0$.

For (ii), parameterize each component by one translation coordinate and its
bounded relative coordinates.  If the close pair lies in one component, it
cuts out a band of width $O(\eta)$ in a bounded relative-coordinate domain.
The logarithmic lattice error is $O(L_n^{-1})$ by
\eqref{eq:harmonic-interval-bound}.  Summing the remaining coordinates gives
$O_{L,M}((\eta+L_n^{-1})T_n^m)$.  If the close pair lies in two different
components, it links two translation coordinates within a bounded interval
whose length depends only on $M$, and therefore removes one free translation
factor.  Its contribution is $O_{L,M}(T_n^{m-1})$.  This proves
\eqref{eq:close-cycle-configurations}.

For (iii), first consider collections with no vertex above $U_n$ but with a
vertex below $L_n$.  Use the maximum-oscillation envelope as in (i).  By the
observation preceding the lemma, the component meeting the lower range has
translation mass $O(1+\log L_n+\Omega)$ rather than $T_n$.  Hence the total is
bounded by
\[
        C_LT_n^{m-1}
        \sum_{k\ge0}(1+\log L_n+k)(2+k)^{L-m}\Psi_L(k)
        =O_L((1+\log L_n)T_n^{m-1})
        =o(T_n^m).
\tag{10.55}\label{eq:lower-boundary-cycle-sum}
\]
The case in which all components are fixed points follows directly from
$\sum_{a<L_n}a^{-1}=O(\log\log n)$ and the same conclusion.

Now suppose that some vertex exceeds $U_n$.  A component lying wholly above
$U_n$ has harmonic mass at most $H_n^r=o(T_n)$ for its fixed length $r$.
For the remaining components, if at least one nontrivial component meets
$[1,U_n]$, choose the one of maximal capped oscillation and use
\eqref{eq:cycle-envelope-maximum}; the summation in (i), together with
\eqref{eq:lower-boundary-cycle-sum}, is then $O_L(T_n^{m-1})$.  If no such
component exists, every remaining component is either a fixed point or wholly
above $U_n$, and its harmonic mass is $O(T_n)$.  Thus any collection
containing a wholly upper component contributes $o(T_n^m)$.

It remains to treat crossing components.  Let $h\ge1$ be the total number of
vertices above $U_n$.  For every component meeting $[1,U_n]$, use the capped
oscillation \eqref{eq:capped-log-oscillation}, and let $\Omega$ be their
maximum.  Corollary \ref{cor:endpoint-cycle-envelope} supplies the factor
$\Psi_L(\Omega)$.  A crossing component has its lower anchor in
$[U_ne^{-\Omega},U_n]$, so it has no free $T_n$ translation factor; its upper
vertices contribute $H_n^h$.  The harmonic-volume bound therefore gives
\[
\begin{aligned}
&\sum_{\Gamma:\ \Gamma\text{ has a crossing component}}
\Pp(E_\Gamma)\\
&\qquad\le C_LH_n^hT_n^{m-1}
\sum_{k\ge0}(2+k)^{L-m+1}\Psi_L(k)
=O_L(H_n^hT_n^{m-1})=o(T_n^m),
\end{aligned}
\tag{10.56}\label{eq:crossing-boundary-cycle-sum}
\]
because $h$ and $L$ are fixed and $(\log\log n)^h=o(\log n)$.  This completes
(iii).
\end{proof}

\begin{proof}[Proof of Theorem \ref{thm:endpoint-short-cycles}]
Let
\[
        m=\sum_{r=1}^Rm_r,
        \qquad
        L=\sum_{r=1}^Rrm_r.
\]
If $m=0$, the assertion is trivial.  Expanding the mixed factorial moment gives
a sum over ordered collections of distinct directed cycles.  Collections
sharing a vertex have probability zero, so only vertex-disjoint collections
contribute.  The falling factorial orders cycles of the same length; hence no
$m_r!$ divisor appears.

Fix $M<\infty$ and $\eta>0$.  Call a collection good if all vertices lie in
$[L_n,U_n]$, every component has logarithmic oscillation at most $M$, and all
$L$ vertices are pairwise separated by more than $\eta$ on logarithmic scale.
For a good collection, every source-to-target ratio lies in
$[e^{-M},e^M]$, and the sorted targets are separated by ratios at least
$e^\eta$.  Proposition \ref{prop:endpoint-general-entry} therefore gives,
uniformly,
\[
        \Pp(E_\Gamma)
        =(1+o(1))\prod_{a\mapsto b\in\Gamma}
        \frac1b k_\gamma\!\left(\log\frac ba\right).
\tag{10.57}\label{eq:good-cycle-entry-product}
\]
The total leading kernel mass of such collections is $O_L(T_n^m)$ by Lemma
\ref{lem:discrete-endpoint-trace} and the nonnegativity of the kernel, so the
accumulated uniform relative error in
\eqref{eq:good-cycle-entry-product} is $o(T_n^m)$.

For one directed cycle of length $r$, Lemma
\ref{lem:discrete-endpoint-trace} gives
\[
        k_\gamma^{*r}(0)T_n+o(T_n)
\]
when a starting vertex is distinguished.  Division by $r$ removes cyclic
rotations.  For an ordered collection of components, the unrestricted leading
kernel sum factors over components.  Equality restrictions and cross-component
close pairs remove $O_L(T_n^{m-1})$, while same-component $\eta$-bands remove
$O_{L,M}(\eta T_n^m)$.  This follows from the harmonic-coordinate proof of
Lemma \ref{lem:negligible-endpoint-cycle-configurations}; on the bounded
oscillation region the leading kernel is bounded, so the same volume estimates
apply verbatim.  Hence the good
contribution is
\[
        \prod_{r=1}^R
        \left(\frac{k_\gamma^{*r}(0)}rT_n\right)^{m_r}
        +o(T_n^m)+O_{L,M}(\eta T_n^m).
\tag{10.58}\label{eq:good-cycle-factorial-main}
\]
The same lemma controls all nongood collections.  Let first $n\to\infty$, then
$\eta\downarrow0$, and finally $M\to\infty$.  Since
\[
        T_n=\log(U_n/L_n)
        =\log n-3\log\log n+O(1)\sim\log n,
\]
\eqref{eq:good-cycle-factorial-main} proves the mixed factorial-moment
asymptotic.

For the $L^p$ conclusion, put
\[
        X_{n,r}=\frac{C_r}{\alpha_{\gamma,r}\log n}.
\]
The first two factorial moments give $X_{n,r}\to1$ in $L^2$, hence in
probability.  Given any fixed $p<\infty$, choose an integer $s>p$.  The
factorial-moment theorem and the identities expressing ordinary powers as
finite linear combinations of falling factorials give
\[
        \sup_n\Ee X_{n,r}^s<\infty.
\]
Thus $|X_{n,r}-1|^p$ is uniformly integrable, and convergence in probability
gives $X_{n,r}\to1$ in $L^p$.  Since $R$ is fixed, the same argument applies
to the vector in any norm on $\mathbb R^R$.
\end{proof}

\begin{remark}[Transpositions]
For $r=2$,
\[
\begin{aligned}
        k_\gamma^{*2}(0)
        &=\int_{\mathbb R}k_\gamma(x)k_\gamma(-x)\dd x\\
        &=\gamma K_\gamma^2
        \int_{\mathbb R}e^{-2K_\gamma\cosh t}\dd t
        =2\gamma K_\gamma^2\mathrm K_0(2K_\gamma),
\end{aligned}
\]
where $t=\gamma x$ and
$\mathrm K_0(z)=\int_0^\infty e^{-z\cosh t}\dd t$.  Thus
\[
        \Ee C_2\sim
        \gamma K_\gamma^2\mathrm K_0(2K_\gamma)\log n.
\]
In the Sukhatme case this coefficient is $\mathrm K_0(2)$.
\end{remark}

% END ARCHIVED BLOCK: endpoint-short-cycles

\begin{thebibliography}{99}

\bibitem{Petrov1975}
V. V. Petrov.
\newblock \emph{Sums of Independent Random Variables}.
\newblock Ergebnisse der Mathematik und ihrer Grenzgebiete, Band 82.
\newblock Springer-Verlag, Berlin--Heidelberg, 1975.

\bibitem{BGT1987}
N. H. Bingham, C. M. Goldie, and J. L. Teugels.
\newblock \emph{Regular Variation}.
\newblock Encyclopedia of Mathematics and its Applications, vol. 27.
\newblock Cambridge University Press, Cambridge, 1987.

\end{thebibliography}

\end{document}
